\chapter{堆、Top K 和贪心证明}

堆和贪心经常一起出现，但它们不是同一个概念。堆是一种数据结构，负责快速取出当前最大或最小候选；贪心是一种算法策略，声称每一步做局部最优选择不会破坏全局最优。堆题的关键是“当前最该处理的候选是谁”；贪心题的关键是“为什么现在做这个选择以后不会后悔”。面试中只写出 \code{priority_queue} 远远不够，还要讲清候选含义、过期规则、比较器方向和正确性理由。

本章先讲堆的候选维护和比较器方向，再讲 Top K、数据流中位数、多路归并；之后再讲贪心证明、跳跃游戏、任务调度、加油站和反悔模型。堆题卡和贪心题卡现在都放在本章末尾，和 DP 章节彻底分开。

C++ 的 \code{std::priority_queue} 默认是大顶堆。要写小顶堆，需要使用 \code{std::greater<T>}。如果元素是结构体，需要自定义比较器。比较器的方向经常让人困惑：\code{priority_queue} 会把“优先级最大”的元素放在 \code{top}，比较器返回 \code{true} 表示左侧优先级低于右侧。因此写自定义比较器时，最好用“谁应该排后面”来理解。

\begin{lstlisting}[language=C++,caption={priority\_queue 的比较器方向}]
struct Entry {
    int value = 0;
    int frequency = 0;
};

struct LowerFrequencyFirst {
    bool operator()(const Entry& lhs, const Entry& rhs) const
    {
        return lhs.frequency > rhs.frequency;
    }
};

using MinHeap = std::priority_queue<Entry,
                                    std::vector<Entry>,
                                    LowerFrequencyFirst>;
\end{lstlisting}

这段比较器让频次更小的元素优先出现在堆顶，因为当 \code{lhs.frequency > rhs.frequency} 时，\code{lhs} 被认为优先级更低。很多人会在这里写反，导致小顶堆变成大顶堆。简单类型能用 \code{std::greater<int>} 时就用标准库；复杂类型必须写清比较器语义，并用小样例验证堆顶是谁。

LeetCode 347 前 K 个高频元素是哈希表和堆的组合题。题目给定整数数组 \code{nums} 和整数 \code{k}，要求返回出现频率最高的 \code{k} 个元素；样例 \code{nums = [1,1,1,2,2,3], k = 2}，输出可以是 \code{[1,2]}。暴力基线是先统计频次，再把所有不同元素按频次完整排序，复杂度 \complexity{O(m\log m)}，\code{m} 是不同元素个数。问题只要前 K 个，不需要完整排名，于是维护一个大小为 K 的小顶堆。堆顶是当前前 K 高频候选里频次最低的元素；新元素频次如果不超过堆顶，就不可能进入前 K；如果超过，就替换堆顶。

\begin{lstlisting}[language=C++,caption={LeetCode 347 前 K 个高频元素：频次表加大小为 K 的小顶堆}]
std::vector<int> top_k_frequent(const std::vector<int>& nums, int k)
{
    std::unordered_map<int, int> frequency;
    frequency.reserve(nums.size());
    for (const int x : nums) {
        ++frequency[x];
    }

    using Entry = std::pair<int, int>;  // frequency, value
    std::priority_queue<Entry, std::vector<Entry>, std::greater<Entry>> heap;

    for (const auto& [value, count] : frequency) {
        heap.emplace(count, value);
        if (static_cast<int>(heap.size()) > k) {
            heap.pop();
        }
    }

    std::vector<int> answer;
    answer.reserve(k);
    while (!heap.empty()) {
        answer.push_back(heap.top().second);
        heap.pop();
    }
    return answer;
}
\end{lstlisting}

这里用 \code{pair<count, value>} 是为了让 \code{std::greater} 按频次小顶堆工作；如果频次相同，会按值比较，但题目通常不要求固定顺序。复杂度是统计 \complexity{O(n)}，堆维护 \complexity{O(m\log k)}。如果 K 接近 m，完整排序可能更简单；如果值域有限，也可以用桶排序按频次分桶，达到线性复杂度。面试时要比较这三种方案，而不是只给一种。

LeetCode 295 数据流的中位数展示双堆模型。题目要求设计 \code{MedianFinder}，支持不断插入数字，并随时返回当前所有数字的中位数；样例操作 \code{addNum(1), addNum(2), findMedian(), addNum(3), findMedian()}，两次查询分别返回 \code{1.5} 和 \code{2.0}。暴力方法每次插入后重新排序，插入代价高；或者维护有序数组，查询容易但插入要搬移元素。中位数只关心较小一半的最大值和较大一半的最小值，于是用一个大顶堆保存较小一半，一个小顶堆保存较大一半。两个堆的大小差不超过 1，并且大顶堆堆顶不大于小顶堆堆顶。这样总数为奇数时，中位数是元素更多那一侧的堆顶；总数为偶数时，中位数是两个堆顶平均值。

\begin{lstlisting}[language=C++,caption={LeetCode 295 数据流中位数：双堆维护两半}]
class MedianFinder {
public:
    void add_num(int num)
    {
        if (this->lower_.empty() || num <= this->lower_.top()) {
            this->lower_.push(num);
        } else {
            this->upper_.push(num);
        }
        this->rebalance();
    }

    double find_median() const
    {
        if (this->lower_.size() == this->upper_.size()) {
            return (static_cast<double>(this->lower_.top()) +
                    static_cast<double>(this->upper_.top())) / 2.0;
        }
        return this->lower_.size() > this->upper_.size()
                   ? static_cast<double>(this->lower_.top())
                   : static_cast<double>(this->upper_.top());
    }

private:
    void rebalance()
    {
        if (this->lower_.size() > this->upper_.size() + 1) {
            this->upper_.push(this->lower_.top());
            this->lower_.pop();
        } else if (this->upper_.size() > this->lower_.size() + 1) {
            this->lower_.push(this->upper_.top());
            this->upper_.pop();
        }
    }

    std::priority_queue<int> lower_;
    std::priority_queue<int, std::vector<int>, std::greater<int>> upper_;
};
\end{lstlisting}

这段代码把重平衡逻辑拆成私有函数，避免 \code{add_num} 过长。成员访问使用 \code{this->}，符合 C++ 代码规范。复杂度是插入 \complexity{O(\log n)}，查询 \complexity{O(1)}。易错点有两个：偶数个元素时要转成 \code{double} 避免整数溢出和整数除法；重平衡只保证大小差，不要忘记插入时先按堆顶决定放入哪一侧，否则大小平衡但顺序不平衡。

LeetCode 23 合并 K 个升序链表是堆在链表上的应用。题目给定 \code{k} 个升序链表，要求合并成一个升序链表；样例 \code{[[1,4,5],[1,3,4],[2,6]]}，输出 \code{[1,1,2,3,4,4,5,6]}。暴力方法可以每次扫描 K 个链表头，找最小节点接入结果，若总节点数为 N，复杂度 \complexity{O(NK)}。两条链表合并也可以两两归并，但若每轮都把一个新链表并入长链，仍可能重复移动很多节点。堆优化是把每条链表当前头节点放入小顶堆，每次弹出最小节点，再把它的后继放入堆。堆里最多 K 个节点，所以复杂度 \complexity{O(N\log K)}。

\begin{lstlisting}[language=C++,caption={LeetCode 23 合并 K 个升序链表：堆保存每条链当前头节点}]
ListNode* merge_k_lists(std::vector<ListNode*> lists)
{
    struct CompareNode {
        bool operator()(const ListNode* lhs, const ListNode* rhs) const
        {
            return lhs->value > rhs->value;
        }
    };

    std::priority_queue<ListNode*, std::vector<ListNode*>, CompareNode> heap;
    for (ListNode* node : lists) {
        if (node != nullptr) {
            heap.push(node);
        }
    }

    ListNode dummy;
    ListNode* tail = &dummy;
    while (!heap.empty()) {
        ListNode* node = heap.top();
        heap.pop();
        if (node->next != nullptr) {
            heap.push(node->next);
        }
        tail->next = node;
        tail = tail->next;
    }
    tail->next = nullptr;
    return dummy.next;
}
\end{lstlisting}

这里比较器把值更大的节点放到后面，因此堆顶是值最小节点。弹出节点后先把后继放入堆，再接入结果，两个顺序都可以，但最后要确保 \code{tail->next = nullptr}，避免结果尾部还连着旧链表的残余路径产生误解。这个算法没有创建新节点，复用输入节点；如果不允许修改输入，需要新建节点。面试表达要强调堆里保存的是“每条链表尚未合并部分的最小候选”。

LeetCode 55 跳跃游戏是贪心的入门题。题目给定数组 \code{nums}，\code{nums[i]} 表示从下标 \code{i} 最多能向右跳多少步，要求判断能否从下标 0 到达最后一个下标；样例 \code{nums = [2,3,1,1,4]} 返回 \code{true}，因为可以从 0 跳到 1，再跳到 4。暴力搜索会从每个位置尝试所有可跳长度，存在大量重复状态。DP 可以记录每个位置是否可达，复杂度 \complexity{O(n^2)}。贪心更简单：扫描过程中维护当前能到达的最远下标 \code{farthest}。如果扫描到的位置已经超过 \code{farthest}，说明这个位置不可达；否则用它继续更新最远范围。

\begin{lstlisting}[language=C++,caption={LeetCode 55 跳跃游戏：维护当前可达最远位置}]
bool can_jump(const std::vector<int>& nums)
{
    int farthest = 0;
    for (int i = 0; i < static_cast<int>(nums.size()); ++i) {
        if (i > farthest) {
            return false;
        }
        farthest = std::max(farthest, i + nums[i]);
        if (farthest >= static_cast<int>(nums.size()) - 1) {
            return true;
        }
    }
    return true;
}
\end{lstlisting}

这个贪心的正确性来自覆盖区间：只要当前位置 \code{i} 可达，那么从它出发能扩展的所有位置都被 \code{farthest} 覆盖；我们不关心具体从哪条路径到达，只关心可达区域的右边界。每次扫描到边界内的新位置，都可能把边界向右推。若某个位置在边界之外，任何之前位置都无法到达它，后面的位置更不可能被访问。

LeetCode 45 跳跃游戏 II 在同样输入上要求到达最后下标的最少跳跃次数；样例 \code{nums = [2,3,1,1,4]}，输出 2，对应路径 \code{0 -> 1 -> 4}。暴力 BFS 可以把每个下标当节点，每条可跳边连接到后续位置，求无权最短路，但边数最坏 \complexity{O(n^2)}。贪心做法把一跳能到达的范围看成 BFS 的一层：\code{current_end} 是当前跳数能覆盖的最右位置，扫描这一层内所有点，计算下一跳能到达的最远位置 \code{farthest}；当扫描到 \code{current_end}，说明这一层结束，必须增加一次跳跃。

\begin{lstlisting}[language=C++,caption={LeetCode 45 跳跃游戏 II：按层扩展可达区间}]
int jump(const std::vector<int>& nums)
{
    if (nums.size() <= 1) {
        return 0;
    }

    int jumps = 0;
    int current_end = 0;
    int farthest = 0;

    for (int i = 0; i < static_cast<int>(nums.size()) - 1; ++i) {
        farthest = std::max(farthest, i + nums[i]);
        if (i == current_end) {
            ++jumps;
            current_end = farthest;
        }
    }

    return jumps;
}
\end{lstlisting}

这段代码本质上是压缩版 BFS。每次增加 \code{jumps}，代表进入下一层；\code{current_end} 以内的位置都属于当前层。为什么循环只到倒数第二个位置？因为一旦已经能到达最后，就不需要从最后一个位置再跳。常见错误是在每个位置都加跳数，或者没有区分当前层边界和下一层最远边界。

贪心证明通常有两种表达：交换论证和领先性论证。区间调度用交换论证：把最优解中结束更晚的选择换成结束更早的选择，不会变差。跳跃游戏用领先性论证：贪心维护的可达最远边界始终不落后于任何其他方案在相同跳数下能到达的边界。面试中给出哪种证明不重要，重要的是不要把“看起来对”当成证明。

LeetCode 621 任务调度器是堆、计数和构造思维的交叉题。题目给定一组任务和冷却时间 \code{n}，相同任务之间必须至少间隔 \code{n} 个单位时间，问完成所有任务的最短时间；样例 \code{tasks = ["A","A","A","B","B","B"], n = 2}，输出 8，一种安排是 \code{A B idle A B idle A B}。暴力模拟可以每个时间点扫描所有任务，选择当前可执行且剩余次数最多的任务；如果任务种类多，扫描会重复。堆优化是把可执行任务按剩余次数放进大顶堆，把冷却中的任务放进队列，时间推进时把冷却结束的任务放回堆。

\begin{lstlisting}[language=C++,caption={LeetCode 621 任务调度器：堆处理可执行任务，队列处理冷却}]
int least_interval(const std::vector<char>& tasks, int n)
{
    std::array<int, 26> count{};
    for (const char task : tasks) {
        ++count[task - 'A'];
    }

    std::priority_queue<int> available;
    for (const int c : count) {
        if (c > 0) {
            available.push(c);
        }
    }

    std::queue<std::pair<int, int>> cooling;  // ready_time, remaining
    int time = 0;
    while (!available.empty() || !cooling.empty()) {
        ++time;
        if (!available.empty()) {
            int remaining = available.top();
            available.pop();
            --remaining;
            if (remaining > 0) {
                cooling.emplace(time + n, remaining);
            }
        }

        if (!cooling.empty() && cooling.front().first == time) {
            available.push(cooling.front().second);
            cooling.pop();
        }
    }

    return time;
}
\end{lstlisting}

这个模拟版本容易解释，但还有数学构造解。设最高频任务出现 \code{max_count} 次，有 \code{max_kinds} 种任务都达到最高频。最短时间至少是 \code{(max_count - 1) * (n + 1) + max_kinds}，因为最高频任务形成 \code{max_count - 1} 个完整间隔块，最后一组放所有最高频任务。答案还不能小于任务总数，所以取二者最大。模拟适合理解过程，公式适合直接求最短长度。面试中给出一种后，最好能说明另一种的直觉。

LeetCode 134 加油站是贪心的区间累积题。题目给定两个数组 \code{gas} 和 \code{cost}，\code{gas[i]} 是第 \code{i} 个站能加到的油，\code{cost[i]} 是从第 \code{i} 个站开到下一个站消耗的油，要求返回能绕一圈的起点下标，若不存在返回 \code{-1}；样例 \code{gas = [1,2,3,4,5], cost = [3,4,5,1,2]}，输出 3。暴力方法从每个站出发模拟一圈，复杂度 \complexity{O(n^2)}。贪心依据是：如果从 \code{start} 出发到 \code{i} 之间某处油量变成负数，那么 \code{start} 到 \code{i} 之间任何站都不可能作为合法起点。因为从这些中间站出发，只会少掉 \code{start} 到该站之前积累的油量，不会更好。

\begin{lstlisting}[language=C++,caption={LeetCode 134 加油站：失败区间整体跳过}]
int can_complete_circuit(const std::vector<int>& gas,
                         const std::vector<int>& cost)
{
    int total = 0;
    int tank = 0;
    int start = 0;

    for (int i = 0; i < static_cast<int>(gas.size()); ++i) {
        const int delta = gas[i] - cost[i];
        total += delta;
        tank += delta;
        if (tank < 0) {
            start = i + 1;
            tank = 0;
        }
    }

    return total >= 0 ? start : -1;
}
\end{lstlisting}

\code{total < 0} 表示全局油量不足，任何起点都不可能成功。若全局油量足够，局部失败时把起点跳到 \code{i + 1} 是安全的。这个跳跃是贪心正确性的核心，不是经验优化。加油站题的复盘要把“失败区间整体排除”说清楚，否则代码看起来像玄学。

LeetCode 763 划分字母区间展示“最后出现位置”如何制造贪心边界。题目要求把字符串划分成尽可能多的片段，使每个字母最多出现在一个片段中，返回每个片段长度；样例 \code{s = "ababcbacadefegdehijhklij"}，输出 \code{[9,7,8]}。暴力方法枚举所有切分并检查字符集合是否交叉，复杂度高。优化思路是先记录每个字符最后出现的位置。扫描时，当前片段右边界必须至少覆盖已经遇到的所有字符的最后位置；当扫描下标到达这个右边界时，当前片段可以安全切开。

\begin{lstlisting}[language=C++,caption={LeetCode 763 划分字母区间：扫描当前片段最远边界}]
std::vector<int> partition_labels(const std::string& s)
{
    std::array<int, 26> last{};
    for (int i = 0; i < static_cast<int>(s.size()); ++i) {
        last[s[i] - 'a'] = i;
    }

    std::vector<int> answer;
    int start = 0;
    int end = 0;
    for (int i = 0; i < static_cast<int>(s.size()); ++i) {
        end = std::max(end, last[s[i] - 'a']);
        if (i == end) {
            answer.push_back(end - start + 1);
            start = i + 1;
        }
    }
    return answer;
}
\end{lstlisting}

这个贪心的边界语义非常清楚：当前片段中出现过的所有字符，它们的最后出现位置都不超过 \code{end}。当 \code{i == end}，说明这些字符不会在后面再次出现，因此切开不会违反题意；如果提前切开，至少有一个字符会跨片段。类似思想也出现在字符串覆盖、最小区间和扫描线问题中：先知道每个对象的最后影响范围，再扫描形成最小合法边界。

LeetCode 502 IPO 项目选择题把贪心和双堆结合起来。题目给定初始资本 \code{w}、最多做 \code{k} 个项目、每个项目需要的启动资本 \code{capital[i]} 和完成后获得的利润 \code{profits[i]}，要求最大化最终资本；样例 \code{k = 2, w = 0, profits = [1,2,3], capital = [0,1,1]}，输出 4。暴力方法每轮扫描所有项目，找当前资本能做且利润最高的项目，复杂度 \complexity{O(kn)}。优化方法是先按资本要求排序，用指针把当前资本可负担的项目加入利润大顶堆，每轮从堆顶选择利润最高项目。

\begin{lstlisting}[language=C++,caption={LeetCode 502 IPO：按资本解锁项目，用利润堆选择}]
int find_maximized_capital(int k,
                           int capital,
                           std::vector<int> profits,
                           std::vector<int> required)
{
    std::vector<std::pair<int, int>> projects;
    projects.reserve(profits.size());
    for (int i = 0; i < static_cast<int>(profits.size()); ++i) {
        projects.emplace_back(required[i], profits[i]);
    }
    std::sort(projects.begin(), projects.end());

    std::priority_queue<int> available_profits;
    int index = 0;
    for (int round = 0; round < k; ++round) {
        while (index < static_cast<int>(projects.size()) &&
               projects[index].first <= capital) {
            available_profits.push(projects[index].second);
            ++index;
        }
        if (available_profits.empty()) {
            break;
        }
        capital += available_profits.top();
        available_profits.pop();
    }

    return capital;
}
\end{lstlisting}

这个问题的局部最优选择是：在当前能做的项目中，选择利润最高的。为什么不会后悔？因为做项目只会增加资本，不会消耗资本；利润越高，下一轮可选项目集合只会更大，不会更小。这个性质非常关键。如果项目会消耗资源，或者项目之间有依赖，简单贪心就可能失败。堆这里只是让“当前可选项目中利润最高”查询更快。

贪心失败的例子同样重要。零钱兑换若面额是 \code{1, 3, 4}，目标 6，贪心先选 4，再选 1、1，共 3 枚；最优是 3、3，共 2 枚。这个反例说明“每次选最大面额”不一定全局最优。判断贪心是否成立，不能只看样例，要找交换论证、单调边界、领先性或 matroid 这类结构。没有证明时，优先退回 DP 或搜索。

LeetCode 767 重构字符串是堆处理“当前最紧急约束”的例子。题目给定字符串 \code{s}，要求重新排列字符，使相邻字符不同；若无法做到，返回空串。样例 \code{s = "aab"} 可以返回 \code{"aba"}，而 \code{"aaab"} 无法重构。暴力回溯可以尝试所有排列，但重复字符多时搜索空间巨大。贪心想法是每次优先使用剩余次数最多的字符，因为它最难安排；但不能连续使用同一个字符，所以每轮从大顶堆取出两个频次最高的不同字符，依次放入答案，再把剩余频次放回堆。若某个字符出现次数超过 \code{(n+1)/2}，必然无法重构。

\begin{lstlisting}[language=C++,caption={LeetCode 767 重构字符串：每次取两个最高频字符}]
std::string reorganize_string(const std::string& s)
{
    std::array<int, 26> count{};
    for (const char ch : s) {
        ++count[ch - 'a'];
    }

    using Entry = std::pair<int, char>;
    std::priority_queue<Entry> heap;
    for (int i = 0; i < static_cast<int>(count.size()); ++i) {
        if (count[i] > 0) {
            heap.emplace(count[i], static_cast<char>('a' + i));
        }
    }

    std::string answer;
    answer.reserve(s.size());
    while (heap.size() >= 2) {
        auto [first_count, first_char] = heap.top();
        heap.pop();
        auto [second_count, second_char] = heap.top();
        heap.pop();

        answer.push_back(first_char);
        answer.push_back(second_char);
        if (--first_count > 0) {
            heap.emplace(first_count, first_char);
        }
        if (--second_count > 0) {
            heap.emplace(second_count, second_char);
        }
    }

    if (!heap.empty()) {
        auto [remaining_count, remaining_char] = heap.top();
        if (remaining_count > 1) {
            return "";
        }
        if (!answer.empty() && answer.back() == remaining_char) {
            return "";
        }
        answer.push_back(remaining_char);
    }

    return answer;
}
\end{lstlisting}

这个贪心的直觉是“高频字符最危险，越晚处理越可能无处可放”。拿 \code{"aaabbc"} 看，\code{a} 有 3 个，其他字符一共 3 个，可以把其他字符形成的空隙想成 \code{_ b _ b _ c _}，最多能放 4 个 \code{a}，所以 3 个 \code{a} 有希望；如果是 \code{"aaaabc"}，\code{a} 有 4 个，其他字符只有 2 个，空隙最多 3 个，必然有相邻 \code{a}。每次取两个不同字符，可以保证刚放入的两个字符相邻不同，也避免同一个最高频字符连续出现。更简洁的写法是每次取一个字符，并把上一轮用过但还剩次数的字符延迟一轮再放回堆；两种写法本质相同。正确性还可以用必要条件解释：最高频字符需要插入到其他字符形成的空隙里，空隙数量最多是其他字符数加一，因此最高频不能超过 \code{(n+1)/2}。

任务调度器和重构字符串类似，但它要求相同任务之间至少间隔 \code{n} 个单位时间，可以插入空闲。模拟做法是每轮取最多 \code{n+1} 个不同任务执行，剩余任务放回堆；如果任务还没做完而本轮不足 \code{n+1}，就需要补空闲。数学做法更快：设最高频为 \code{maxFreq}，拥有最高频的任务种类数为 \code{maxCount}，最短时间至少是 \code{(maxFreq - 1) * (n + 1) + maxCount}，再和任务总数取最大。

\begin{lstlisting}[language=C++,caption={LeetCode 621 任务调度器：按最高频任务计算骨架长度}]
int least_interval(const std::vector<char>& tasks, int cooldown)
{
    std::array<int, 26> count{};
    for (const char task : tasks) {
        ++count[task - 'A'];
    }

    int max_frequency = 0;
    int max_count = 0;
    for (const int frequency : count) {
        if (frequency > max_frequency) {
            max_frequency = frequency;
            max_count = 1;
        } else if (frequency == max_frequency) {
            ++max_count;
        }
    }

    const int skeleton =
        (max_frequency - 1) * (cooldown + 1) + max_count;
    return std::max(static_cast<int>(tasks.size()), skeleton);
}
\end{lstlisting}

公式的含义是：最高频任务形成 \code{max_frequency - 1} 个完整间隔块，每块长度至少 \code{cooldown + 1}，最后再放所有最高频任务的尾部。若其他任务足够多，它们可以填满空隙，答案就是任务总数；若不够，就必须空闲，答案由骨架决定。这个题很好地区分了“堆模拟”和“数学贪心”：堆模拟更通用，公式更简洁但依赖对结构的深理解。

LeetCode 630 课程表 III 是“按截止时间排序 + 最大堆反悔”的经典题。题目给定若干课程 \code{[duration,lastDay]}，每门课需要连续学习 \code{duration} 天，必须不晚于 \code{lastDay} 完成，要求最多能上多少门；样例 \code{courses = [[100,200],[200,1300],[1000,1250],[2000,3200]]}，输出 3。暴力搜索选课集合是指数级。贪心先按截止时间升序处理课程，尽量把当前课程加入计划；如果加入后总时间超过当前截止时间，就从已选课程中删除耗时最长的一门。为什么删最长？因为课程数量减少一门时，删除耗时最长能让总时间降得最多，给后续课程留下最大空间。

\begin{lstlisting}[language=C++,caption={LeetCode 630 课程表 III：超过截止时删除最长课程}]
int schedule_course(std::vector<std::vector<int>> courses)
{
    std::sort(courses.begin(), courses.end(),
              [](const std::vector<int>& lhs, const std::vector<int>& rhs) {
                  return lhs[1] < rhs[1];
              });

    std::priority_queue<int> selected_durations;
    int total_time = 0;
    for (const std::vector<int>& course : courses) {
        const int duration = course[0];
        const int deadline = course[1];
        selected_durations.push(duration);
        total_time += duration;

        if (total_time > deadline) {
            total_time -= selected_durations.top();
            selected_durations.pop();
        }
    }

    return static_cast<int>(selected_durations.size());
}
\end{lstlisting}

这个算法的堆不是为了取下一门课，而是为了在当前已选集合中执行“反悔”。按截止时间排序保证当处理当前课程时，只需要满足当前及更早截止时间；如果超时，必须删掉一门已选课程。删最长课程不会减少已选课程数量以外的可行性，且让总耗时最小。这个证明属于交换论证：如果某个可行方案删的是较短课程而保留较长课程，交换为保留较短、删除较长，课程数量不变，总时间不增。

LeetCode 435 无重叠区间也能从“反悔”角度理解。题目给定若干区间，要求删除最少数量的区间，使剩余区间互不重叠；样例 \code{intervals = [[1,2],[2,3],[3,4],[1,3]]}，输出 1，删除 \code{[1,3]} 即可。暴力枚举删除哪些区间是指数级。按起点排序扫描时，若当前区间和上一个保留区间重叠，就必须删除其中一个。为了给后续留下更多空间，应保留结束更早的那个，删除结束更晚的那个。按终点排序的版本更直接：每次选择结束最早且与已选区间不重叠的区间，最大化保留数量；答案也可以用总数减去最多保留数量。

\begin{lstlisting}[language=C++,caption={LeetCode 435 无重叠区间：保留结束更早的区间}]
int erase_overlap_intervals(std::vector<std::vector<int>> intervals)
{
    if (intervals.empty()) {
        return 0;
    }
    std::sort(intervals.begin(), intervals.end(),
              [](const std::vector<int>& lhs, const std::vector<int>& rhs) {
                  return lhs[0] < rhs[0];
              });

    int removed = 0;
    int current_end = intervals[0][1];
    for (int i = 1; i < static_cast<int>(intervals.size()); ++i) {
        if (intervals[i][0] >= current_end) {
            current_end = intervals[i][1];
            continue;
        }
        ++removed;
        current_end = std::min(current_end, intervals[i][1]);
    }
    return removed;
}
\end{lstlisting}

当发生重叠时，\code{current_end = min(current_end, intervals[i][1])} 表示保留结束更早的区间。这个写法不显式记录保留哪个区间，但保留了对后续唯一重要的信息：当前已保留区间的最小结束时间。贪心题经常不需要保留完整历史，只保留能决定未来可行性的摘要状态。若题目要求输出具体删除哪些区间，就需要额外记录选择。

LeetCode 406 根据身高重建队列是排序和贪心的组合。题目给定若干人 \code{[height,k]}，其中 \code{k} 是排在该人前面且身高大于等于 \code{height} 的人数，要求重建队列；样例 \code{[[7,0],[4,4],[7,1],[5,0],[6,1],[5,2]]}，输出可以是 \code{[[5,0],[7,0],[5,2],[6,1],[4,4],[7,1]]}。暴力可以尝试所有排列并检查每个人的 \code{k}，复杂度不可接受。若从矮到高放，后面更高的人插入会影响矮个子的 \code{k}；若从高到矮放，已经放入的人都不矮于当前人，当前人的 \code{k} 就等于插入位置。排序规则是身高降序，身高相同则 \code{k} 升序，然后按 \code{k} 插入结果数组。

\begin{lstlisting}[language=C++,caption={LeetCode 406 根据身高重建队列：高个先放}]
std::vector<std::vector<int>> reconstruct_queue(
    std::vector<std::vector<int>> people)
{
    std::sort(people.begin(), people.end(),
              [](const std::vector<int>& lhs, const std::vector<int>& rhs) {
                  if (lhs[0] != rhs[0]) {
                      return lhs[0] > rhs[0];
                  }
                  return lhs[1] < rhs[1];
              });

    std::vector<std::vector<int>> queue;
    queue.reserve(people.size());
    for (const std::vector<int>& person : people) {
        queue.insert(queue.begin() + person[1], person);
    }
    return queue;
}
\end{lstlisting}

这段代码的插入是 \complexity{O(n)}，整体 \complexity{O(n^2)}，但 LeetCode 规模通常可接受。若规模很大，可以用树状数组或线段树找第 k 个空位。正确性来自处理顺序：当插入当前人时，队列中已有的人都比他高或等高，因此插入下标正好决定了他前面符合条件的人数；后续更矮的人插入不会影响他的 \code{k}。同高的人按 \code{k} 升序处理，避免同高较大 k 先插入造成位置不稳定。

合并石头、戳气球、最优二叉搜索树这类题看上去也有“每次选局部最小或最大”的诱惑，但通常不能简单贪心。以合并石头为例，如果每次合并当前最小的相邻两堆，不一定得到全局最优，因为相邻约束会改变未来可合并结构。没有相邻约束的哈夫曼合并可以每次取最小两项，这是因为任意两项都可合并，且最优前缀树有明确交换性质；有相邻约束后，问题变成区间 DP。这个对比非常重要：同样是“合并代价”，约束不同，算法类型就不同。

LeetCode 860 柠檬水找零是简单贪心，但很适合训练“优先保留灵活资源”。题目中每杯柠檬水 5 元，顾客按顺序支付 5、10 或 20，要求判断能否为每位顾客正确找零；样例 \code{bills = [5,5,5,10,20]} 返回 \code{true}。暴力思路可以把手里的纸币组合都枚举出来，但这里币值很少，只需要维护 5 元和 10 元数量。收到 10 需要找一个 5；收到 20 时，优先找 10+5，而不是三个 5，因为 10 只能在找 20 时使用，5 更灵活，后续找 10 和 20 都需要。这个证明不是复杂数学，而是资源灵活性：越通用的资源越应该保留。若无法找零则返回 false。

\begin{lstlisting}[language=C++,caption={LeetCode 860 柠檬水找零：优先使用较不灵活的 10 元}]
bool lemonade_change(const std::vector<int>& bills)
{
    int five = 0;
    int ten = 0;
    for (const int bill : bills) {
        if (bill == 5) {
            ++five;
        } else if (bill == 10) {
            if (five == 0) {
                return false;
            }
            --five;
            ++ten;
        } else {
            if (ten > 0 && five > 0) {
                --ten;
                --five;
            } else if (five >= 3) {
                five -= 3;
            } else {
                return false;
            }
        }
    }
    return true;
}
\end{lstlisting}

这个题也能说明贪心证明可以很朴素。不是所有贪心都需要长篇交换论证，有时只要明确资源的支配关系：一张 5 元可以参与两类找零，一张 10 元只能参与 20 元找零；在能完成当前找零的前提下，消耗更不灵活的资源不会让未来更差。若题目币值变化，这个策略可能失效，因此不能把它泛化成“永远优先用大面额”。

LeetCode 871 最低加油次数是堆和贪心的又一个“反悔”模型。题目给定目标距离 \code{target}、初始油量 \code{startFuel} 和若干加油站 \code{[position,fuel]}，要求到达目标的最少加油次数；样例 \code{target = 100, startFuel = 10, stations = [[10,60],[20,30],[30,30],[60,40]]}，输出 2。暴力搜索在哪些站加油是指数级。贪心扫描所有可达加油站，把它们的油量放入大顶堆；当当前油量无法到达下一个站或终点时，就从已经路过且可选择的站中取油量最大的一个加油。虽然看起来像回头加油，但等价于在之前经过时选择了那个站。

\begin{lstlisting}[language=C++,caption={LeetCode 871 最小加油次数：不可达时从已过站点中选最大油量}]
int min_refuel_stops(int target,
                     int start_fuel,
                     std::vector<std::vector<int>> stations)
{
    stations.push_back({target, 0});
    std::priority_queue<int> fuels;
    int fuel = start_fuel;
    int previous_position = 0;
    int stops = 0;

    for (const std::vector<int>& station : stations) {
        const int position = station[0];
        const int available_fuel = station[1];
        fuel -= position - previous_position;
        while (fuel < 0 && !fuels.empty()) {
            fuel += fuels.top();
            fuels.pop();
            ++stops;
        }
        if (fuel < 0) {
            return -1;
        }
        fuels.push(available_fuel);
        previous_position = position;
    }

    return stops;
}
\end{lstlisting}

这个算法的正确性来自延迟决策：在还没被迫加油之前，不决定具体在哪个已过站点加；一旦必须加油，为了让未来最容易可达，就选择已过站点中油量最大的。选择更小油量不会减少加油次数，反而可能更早再次陷入不可达。堆中保存的是已经经过但尚未使用的加油机会。这个模型和 IPO、课程表 III 都有相似结构：先收集可用候选，必要时选最有利的一个。

贪心题的讲解应该包含失败测试。跳跃游戏 II 如果每一步只跳到当前能跳到的最远下标，可能不是最少跳数；正确做法是按当前层覆盖范围推进，选择下一层最远覆盖。课程表 III 如果只按持续时间短排序，会错过截止时间约束；如果只按截止时间排序但不反悔，也会被长课程拖垮。重构字符串如果每次只取最高频字符但不延迟上一字符，就可能产生相邻重复。把这些失败原因讲出来，答案才可信。

\begin{labbox}
堆与贪心实验：第一，为零钱兑换构造 \code{1,3,4} 和目标 6，写出贪心失败过程，再用 DP 修正。第二，给课程表 III 构造一门很长但截止早的课程，观察最大堆如何反悔。第三，用任务调度器分别实现堆模拟和公式法，比较输出一致性。第四，把最小加油次数中的大顶堆改成小顶堆，找一个用例说明为什么会增加加油次数或失败。
\end{labbox}

\begin{principlebox}{用案例复盘堆和贪心}
LeetCode 215 数组第 K 大里，堆顶保存当前前 K 大的边界，和完整排序、快速选择形成对比。LeetCode 295 数据流中位数里，两个堆分别保存较小一半和较大一半，大小和平衡顺序两个不变量缺一不可。LeetCode 23 合并 K 个链表里，堆中每条链只放当前头节点，弹出后推进该链。LeetCode 435 和 452 的贪心选择来自按右端点的交换论证；LeetCode 630 课程表 III 的反悔来自“先接受，再在超时后丢掉持续时间最长课程”。复盘时把题先归到这些案例，再说明堆里保存什么、堆顶为何可用、局部选择为什么不伤害全局。
\end{principlebox}

堆和动态极值深度题卡 按题型拆成章内大节，但每个大节都先从 LeetCode 案例切入，再进入分流、案例矩阵、案例推导和实验复盘。这样读者看到的是从具体题推到规律，而不是先读抽象方法再猜它能用在哪。

\topic{堆与动态极值}
这一节先看 LeetCode 23 合并 K 个升序链表、LeetCode 215 数组中的第 K 个最大元素、LeetCode 295 数据流的中位数 这几道 LeetCode 题，再整理同一题型的分流和推导。阅读顺序不是先背原理，而是先看案例的暴力瓶颈，再把重复出现的观察抽象成方法。

\subsection{原理和适用条件}

以 LeetCode 23 合并 K 个升序链表、LeetCode 215 数组中的第 K 个最大元素、LeetCode 295 数据流的中位数 为主线：LeetCode 23 合并 K 个升序链表：模型是“每条链表当前头节点都是可能的下一个最小元素。”，优化抓手是“小顶堆保存每条非空链的当前头，弹出最小节点后把它的后继入堆。”；LeetCode 215 数组中的第 K 个最大元素：模型是“只关心第 K 大，不需要完整排序。”，优化抓手是“大小为 K 的小顶堆保存当前前 K 大边界；快速选择可做到平均线性。”；LeetCode 295 数据流的中位数：模型是“数据流需要动态维护较小一半和较大一半。”，优化抓手是“大顶堆保存较小一半，小顶堆保存较大一半，大小差不超过一。”。这些案例共同推出的规律是：先每轮扫描所有候选找极值，再把动态极值查询交给堆。堆只保证堆顶，不保证全局遍历有序；有过期元素时要在使用堆顶前清理。 方法名只是复盘后的标签，真正要掌握的是从题面约束、暴力失败、状态设计到边界反例的推导链。

\begin{principlebox}{图示：从 LeetCode 案例到 堆与动态极值推导链}
\begin{tabular}{@{}p{0.18\linewidth}p{0.74\linewidth}@{}}
暴力基线 & 枚举候选、完整模拟或逐个检查，先保证覆盖全部可能答案。\\
瓶颈定位 & 慢点在于动态极值查询。我们只需要当前最极端元素，不需要整个集合完全有序。\\
结构选择 & 堆把插入和弹出极值控制在对数复杂度。Top K 用固定大小堆维护边界，合并多路序列用堆保存每一路当前头部，数据流中位数用双堆维护两半。\\
不变量 & 堆不变量不是全局有序，而是堆顶代表当前最优候选。若存在过期元素，需要在弹出时懒删除，保证真正使用堆顶前它仍然有效。\\
代码落地 & C++ 的 \code{std::priority_queue} 默认大顶堆。小顶堆用 \code{std::greater<T>}；结构体比较器要按谁优先级更低来写。堆中保存大对象时尽量保存下标或指针。\\
\end{tabular}
\end{principlebox}

\subsection{LeetCode 案例分流}

\begin{principlebox}{从 LeetCode 案例分流：堆与动态极值}
\textbf{Top K 和动态排名。}
代表案例：LeetCode 215 数组第 K 大、LeetCode 347 前 K 高频、LeetCode 973 最接近原点 K 个点。先看这些题的题面条件：只需要前 K 个、最近 K 个或第 K 个，不需要完整排序。由此推出的处理方式是：维护固定大小堆，堆顶是当前边界。风险点：K 的边界、比较器方向和输出顺序要说明。

\textbf{多路归并。}
代表案例：LeetCode 23 合并 K 个升序链表、LeetCode 632 最小区间覆盖 K 个列表。先看这些题的题面条件：多个有序序列各自只能暴露当前头部。由此推出的处理方式是：堆保存每一路当前候选，弹出后推进该路。风险点：堆里最多每路一个候选，某一路耗尽会改变终止条件。

\textbf{双堆和数据流。}
代表案例：LeetCode 295 数据流中位数、LeetCode 480 滑动窗口中位数。先看这些题的题面条件：数据在线到达，需要动态维护中位数或两侧平衡。由此推出的处理方式是：大顶堆保存较小一半，小顶堆保存较大一半。风险点：若支持删除，需要延迟删除，堆顶使用前必须清理。

\textbf{调度和反悔。}
代表案例：LeetCode 621 任务调度器、LeetCode 630 课程表 III、LeetCode 1675 数组最小偏移量。先看这些题的题面条件：每一步选择当前收益最大、冷却结束或最紧急候选。由此推出的处理方式是：用堆组织当前可选任务或可反悔候选。风险点：要区分堆模拟和数学公式，二者适用条件不同。

\end{principlebox}

\subsection{案例矩阵}

本节案例覆盖 LeetCode 23 合并 K 个升序链表、LeetCode 215 数组中的第 K 个最大元素、LeetCode 295 数据流的中位数、LeetCode 347 前 K 个高频元素、LeetCode 480 滑动窗口中位数、LeetCode 502 IPO、LeetCode 621 任务调度器、LeetCode 632 最小区间覆盖 K 个列表、LeetCode 703 数据流中的第 K 大元素、LeetCode 767 重构字符串、LeetCode 857 雇佣 K 名工人的最低成本、LeetCode 871 最低加油次数、LeetCode 973 最接近原点的 K 个点、LeetCode 1046 最后一块石头的重量、LeetCode 1353 最多可以参加的会议数目、LeetCode 1675 数组的最小偏移量、LeetCode 1851 包含每个查询的最小区间。阅读时不要按题号背答案，而要先判断题面落在哪个分流场景：查询目标是什么，输入约束给了什么单调性或等价关系，暴力慢在哪里，优化结构保存了什么状态。

\begin{principlebox}{案例矩阵}
\textbf{LeetCode 23 合并 K 个升序链表。}
每条链表当前头节点都是可能的下一个最小元素。单点风险：空链表数组、某条链为空、重复值、堆里保存节点指针还是值。

\textbf{LeetCode 215 数组中的第 K 个最大元素。}
只关心第 K 大，不需要完整排序。单点风险：K 为一、K 为 n、重复值、随机 pivot。

\textbf{LeetCode 295 数据流的中位数。}
数据流需要动态维护较小一半和较大一半。单点风险：偶数个元素、负数、重复值、平均值溢出。

\textbf{LeetCode 347 前 K 个高频元素。}
先统计频次，再只维护前 K 个候选。单点风险：K 等于不同元素数、频次相同、结果顺序不要求。

\textbf{LeetCode 480 滑动窗口中位数。}
窗口移动时既要加入新元素，也要删除滑出的旧元素并查询中位数。单点风险：重复值、偶数窗口、平均值溢出、堆顶过期。

\textbf{LeetCode 502 IPO。}
当前资本决定可解锁项目集合，目标是在最多 k 次选择中最大化最终资本。单点风险：初始资本为零、没有可做项目、k 大于项目数、利润为零、资本解锁链式增长。

\textbf{LeetCode 621 任务调度器。}
相同任务之间有冷却间隔，核心受最高频任务骨架限制。单点风险：冷却为零、多个最高频、空闲被其他任务填满。

\textbf{LeetCode 632 最小区间覆盖 K 个列表。}
每个列表至少取一个元素时，当前区间由最小当前值和最大当前值决定。单点风险：某个列表为空、重复值、区间长度相同的 tie-breaker、列表耗尽。

\textbf{LeetCode 703 数据流中的第 K 大元素。}
数据不断到达，只关心当前第 K 大而非完整排序。单点风险：初始元素少于 K、重复值、K 为一、连续 add。

\textbf{LeetCode 767 重构字符串。}
相邻不能相同，剩余次数最多的字符最难安排。单点风险：单字符、最高频过半、多个字符同频、最后只剩一个字符、结果不唯一。

\textbf{LeetCode 857 雇佣 K 名工人的最低成本。}
工资比例由当前队伍中最高工资质量比决定，总成本等于比例乘质量和。单点风险：K 为一、比例相同、浮点精度、质量和更新。

\textbf{LeetCode 871 最低加油次数。}
行驶过程中，经过的加油站形成可反悔的燃料候选。单点风险：起点油足够、无法到达第一个站、多个站同位置、目标作为终点事件。

\textbf{LeetCode 973 最接近原点的 K 个点。}
距离只需比较平方，避免开方误差和成本。单点风险：K 为一、K 为全部、距离相同、坐标平方溢出。

\textbf{LeetCode 1046 最后一块石头的重量。}
每次都要取当前最重的两块石头碰撞。单点风险：没有石头、一块石头、两块相等、多次产生相同重量。

\textbf{LeetCode 1353 最多可以参加的会议数目。}
每天只能参加一个正在进行的会议，应优先选择最早结束的会议。单点风险：同一天多个会议、过期会议清理、空闲日期跳跃、结束日相同。

\textbf{LeetCode 1675 数组的最小偏移量。}
奇数可以先乘二，偶数可以不断除二，目标是缩小当前最大最小差。单点风险：全奇数、全偶数、最大值变奇数停止、数值范围。

\textbf{LeetCode 1851 包含每个查询的最小区间。}
每个查询点需要找覆盖它的最短区间，区间和查询都可离线排序。单点风险：没有覆盖区间、区间长度相同、查询重复、右端过期。

\end{principlebox}

\subsection{共性落点}

\begin{principlebox}{本组共性落点}
\textbf{慢版本。} 暴力做法每轮扫描所有候选找最大或最小，或者每次变化后重新排序。若候选集合动态变化，这会把大量旧排序工作丢掉。
\textbf{瓶颈。} 慢点在于动态极值查询。我们只需要当前最极端元素，不需要整个集合完全有序。
\textbf{状态字段。} 堆元素字段、堆顶比较规则、候选是否过期、答案读取时机；如果需要懒删除，还要保存删除计数。
\textbf{不变量。} 堆不变量不是全局有序，而是堆顶代表当前最优候选。若存在过期元素，需要在弹出时懒删除，保证真正使用堆顶前它仍然有效。
\textbf{实现检查。} 比较器方向先用小样例验证；每次使用堆顶前清理过期项；堆中保存下标时要能回查原数据。
\textbf{C++ 落地。} C++ 的 \code{std::priority_queue} 默认大顶堆。小顶堆用 \code{std::greater<T>}；结构体比较器要按谁优先级更低来写。堆中保存大对象时尽量保存下标或指针。
\textbf{证明抓手。} 证明从候选覆盖开始：所有可能成为下一步答案的对象都在堆中，堆顶过期时不会被使用。
\end{principlebox}

\subsection{案例推导}

\noindent\textbf{LeetCode 23 合并 K 个升序链表。}

\textbf{题目说明。} LeetCode 23 合并 K 个升序链表 的题面入口要先还原成具体任务：每条链表当前头节点都是可能的下一个最小元素。

\textbf{样例入口。} 拿三条链表头 1、1、2 手算，每次答案只需要当前最小头节点；弹出某条链的头后，只有它的后继会成为新候选。暴力每轮扫 k 个头，堆把“找当前最小头”变成对数操作。

\textbf{暴力基线。} 暴力每轮扫描 K 个链表头，找到当前最小头节点接到结果后推进该链。

\textbf{暴力过程。} 若三条链当前头是 1、1、2，第一步只需从这三个头里选最小；弹出某条链的 1 后，只有它的后继会成为新候选。

\textbf{重复工作。} 题面模型：每条链表当前头节点都是可能的下一个最小元素。暴力版的慢点要落到本题：浪费在每输出一个节点都线性扫 K 个头。候选集合动态变化，但每次只需要最小头节点，适合小顶堆。后面的优化只消掉这类重复工作，不能改变答案集合。

\textbf{优化条件。} 本题的关键观察是：小顶堆保存每条非空链的当前头，弹出最小节点后把它的后继入堆。这句话必须能翻译成可维护状态；如果题目条件改变后观察不再成立，就不能继续套原来的写法。

\textbf{相邻状态。} 规律是堆里始终最多保存每条非空链的一个当前头。把这个特例继续拆开看：堆顶必须正好代表下一步最需要处理的候选；若堆里可能有过期对象，读取堆顶前要先清理。堆只解决动态极值，不自动证明被弹出或被丢弃的元素无用。

\textbf{状态复用。} 小顶堆保存每条非空链的当前头节点。弹出最小节点接入结果；若该节点有 next，把 next 入堆。手算时只改被新输入影响的那一小部分，而不是回头重算完整候选。

\textbf{指针和字段。} 状态字段要能说清：heap 中元素是 ListNode*，比较依据是 node->val；tail 是结果尾；node->next 是该链的新候选。手算时按这个协议跟踪：头节点 1、1、2 入堆；弹出第一个 1 后把它的后继入堆，堆仍保存每条链当前最小暴露节点。

\textbf{代码落点。} 关键是堆里保存节点指针而不是复制值，这样可以复用原节点。比较器要写成小顶堆语义，C++ priority\_queue 默认是大顶堆。关键代码行必须能逐句翻译成“旧状态怎样变成新状态”，否则就只是模板。

\textbf{正确性。} 每条有序链未输出部分的最小值就是当前头；所有链当前头中的最小者就是全局下一个输出节点。

\textbf{边界实验。} 先用空链表数组、某条链为空、重复值、堆里保存节点指针还是值做反例检查。实验时覆盖空链表数组、某条链为空、重复值、堆里保存节点指针还是值，记录堆顶是否过期、比较器是否让正确候选在堆顶、K 或窗口大小为边界值时是否稳定。

\textbf{复杂度变化。} 每轮扫 K 个头是 O(NK)。堆中最多 K 个节点，每个节点入堆出堆一次，总时间 O(N log K)，空间 O(K)。

\textbf{迁移追问。} 如果题目改成“也可以分治两两合并，堆更适合 K 较大且链长不均的场景”，先判断原状态是否仍然覆盖新答案集合，再决定是微调边界、改 value 语义，还是换算法。

\noindent\textbf{LeetCode 215 数组中的第 K 个最大元素。}

\textbf{题目说明。} LeetCode 215 数组中的第 K 个最大元素 的题面入口要先还原成具体任务：只关心第 K 大，不需要完整排序。

\textbf{样例入口。} 拿 [3, 2, 1, 5, 6, 4]、k=2 手算，固定大小为 2 的小顶堆最终保存最大的两个数，堆顶就是第 2 大。

\textbf{暴力基线。} 暴力排序整个数组，然后返回降序第 k 个或升序第 n-k 个。

\textbf{暴力过程。} 以 [3, 2, 1, 5, 6, 4]、k=2 为例，完整排序后第 2 大是 5；但为了一个第 K 大，不需要知道所有元素的完整顺序。

\textbf{重复工作。} 题面模型：只关心第 K 大，不需要完整排序。暴力版的慢点要落到本题：浪费在完整排序。我们只需要维护当前最大的 K 个元素，其中最小的那个就是第 K 大边界。后面的优化只消掉这类重复工作，不能改变答案集合。

\textbf{优化条件。} 本题的关键观察是：大小为 K 的小顶堆保存当前前 K 大边界；快速选择可做到平均线性。这句话必须能翻译成可维护状态；如果题目条件改变后观察不再成立，就不能继续套原来的写法。

\textbf{相邻状态。} 规律是堆里只保留当前前 k 大候选，比堆顶小的数不可能进入答案。把这个特例继续拆开看：堆顶必须正好代表下一步最需要处理的候选；若堆里可能有过期对象，读取堆顶前要先清理。堆只解决动态极值，不自动证明被弹出或被丢弃的元素无用。

\textbf{状态复用。} 大小为 K 的小顶堆保存当前前 K 大。新数大于堆顶时，弹出堆顶并加入新数；小于等于堆顶则不可能进入前 K。手算时只改被新输入影响的那一小部分，而不是回头重算完整候选。

\textbf{指针和字段。} 状态字段要能说清：heap.top 是当前前 K 大中的最小值，size 保持不超过 K，num 是当前扫描元素。手算时按这个协议跟踪：扫描样例时，堆最终保存 5 和 6，堆顶 5 就是第 2 大。

\textbf{代码落点。} 关键是使用小顶堆，C++ priority\_queue 默认大顶堆，需要 greater<int>。若选择快速选择，要处理随机 pivot 和重复值。关键代码行必须能逐句翻译成“旧状态怎样变成新状态”，否则就只是模板。

\textbf{正确性。} 堆中始终保存已扫描元素的前 K 大；若新元素不超过堆顶，它最多排在第 K 大之后，可以安全丢弃。

\textbf{边界实验。} 先用K 为一、K 为 n、重复值、随机 pivot做反例检查。实验时覆盖K 为一、K 为 n、重复值、随机 pivot，记录堆顶是否过期、比较器是否让正确候选在堆顶、K 或窗口大小为边界值时是否稳定。

\textbf{复杂度变化。} 排序 O(n log n)，小顶堆 O(n log k) 时间、O(k) 空间；快速选择平均 O(n)。

\textbf{迁移追问。} 如果题目改成“若数据流持续到来，堆方案比一次性快速选择更自然”，先判断原状态是否仍然覆盖新答案集合，再决定是微调边界、改 value 语义，还是换算法。

\noindent\textbf{LeetCode 295 数据流的中位数。}

\textbf{题目说明。} LeetCode 295 数据流的中位数 的题面入口要先还原成具体任务：数据流需要动态维护较小一半和较大一半。

\textbf{样例入口。} 拿数据流 1、2、3 手算，插入 1 后中位数 1；插入 2 后两半是 [1] 和 [2]，中位数 1.5；插入 3 后右半多一个，要平衡成左半 [2, 1]、右半 [3]。

\textbf{暴力基线。} 暴力每插入一个数就重新排序所有已到达元素，然后按长度奇偶取中位数。

\textbf{暴力过程。} 数据流 1、2、3 中，插入 1 中位数是 1；插入 2 后中位数是 (1+2)/2；插入 3 后中位数是 2。

\textbf{重复工作。} 题面模型：数据流需要动态维护较小一半和较大一半。暴力版的慢点要落到本题：浪费在每次都维护完整排序。中位数只需要较小一半的最大值和较大一半的最小值。后面的优化只消掉这类重复工作，不能改变答案集合。

\textbf{优化条件。} 本题的关键观察是：大顶堆保存较小一半，小顶堆保存较大一半，大小差不超过一。这句话必须能翻译成可维护状态；如果题目条件改变后观察不再成立，就不能继续套原来的写法。

\textbf{相邻状态。} 规律是大顶堆保存较小一半，小顶堆保存较大一半，大小差最多 1。把这个特例继续拆开看：堆顶必须正好代表下一步最需要处理的候选；若堆里可能有过期对象，读取堆顶前要先清理。堆只解决动态极值，不自动证明被弹出或被丢弃的元素无用。

\textbf{状态复用。} 大顶堆 small 保存较小一半，小顶堆 large 保存较大一半。保持 small.size() 与 large.size() 差不超过 1，且 small.top()<=large.top()。手算时只改被新输入影响的那一小部分，而不是回头重算完整候选。

\textbf{指针和字段。} 状态字段要能说清：small.top 是左半最大值，large.top 是右半最小值；奇数时较大堆顶是中位数，偶数时两个堆顶平均。手算时按这个协议跟踪：插入 1 放 small；插入 2 后移动平衡成 small=[1]、large=[2]；插入 3 后 large 多，再移动 2 到 small，中位数 2。

\textbf{代码落点。} 关键是先按顺序不变量放入堆，再重平衡大小。求平均时用 double，并防止两个 int 相加溢出。关键代码行必须能逐句翻译成“旧状态怎样变成新状态”，否则就只是模板。

\textbf{正确性。} 两个堆按大小分割数据流，中位数只可能位于左半最大或右半最小；大小不变量保证取堆顶即可。

\textbf{边界实验。} 先用偶数个元素、负数、重复值、平均值溢出做反例检查。实验时覆盖偶数个元素、负数、重复值、平均值溢出，记录堆顶是否过期、比较器是否让正确候选在堆顶、K 或窗口大小为边界值时是否稳定。

\textbf{复杂度变化。} 每次重新排序 O(n log n)；双堆 addNum O(log n)，findMedian O(1)，空间 O(n)。

\textbf{迁移追问。} 如果题目改成“若还要删除任意元素，需要延迟删除哈希表”，先判断原状态是否仍然覆盖新答案集合，再决定是微调边界、改 value 语义，还是换算法。

\noindent\textbf{LeetCode 347 前 K 个高频元素。}

\textbf{题目说明。} LeetCode 347 前 K 个高频元素 的题面入口要先还原成具体任务：先统计频次，再只维护前 K 个候选。

\textbf{样例入口。} 拿 [1, 1, 1, 2, 2, 3]、k=2 手算，频次是 1->3、2->2、3->1，前两个高频是 1 和 2。

\textbf{暴力基线。} 暴力先统计频次，再把所有不同元素按频次完整排序，取前 K 个。

\textbf{暴力过程。} 以 [1, 1, 1, 2, 2, 3]、k=2 为例，频次为 1->3、2->2、3->1，答案是 1 和 2。

\textbf{重复工作。} 题面模型：先统计频次，再只维护前 K 个候选。暴力版的慢点要落到本题：浪费在完整排序所有不同元素。只要前 K 高频，低于当前第 K 高频边界的元素可以丢弃。后面的优化只消掉这类重复工作，不能改变答案集合。

\textbf{优化条件。} 本题的关键观察是：小顶堆大小超过 K 就弹出最低频；桶排序利用频次上界。这句话必须能翻译成可维护状态；如果题目条件改变后观察不再成立，就不能继续套原来的写法。

\textbf{相邻状态。} 规律是先用哈希计数，再用大小为 k 的小顶堆或桶按频次取 Top K；频次相同不影响题目集合答案。把这个特例继续拆开看：堆顶必须正好代表下一步最需要处理的候选；若堆里可能有过期对象，读取堆顶前要先清理。堆只解决动态极值，不自动证明被弹出或被丢弃的元素无用。

\textbf{状态复用。} 先用哈希表计数。再用大小为 K 的小顶堆保存候选，堆顶是当前前 K 中频次最低者；超过 K 就弹出。手算时只改被新输入影响的那一小部分，而不是回头重算完整候选。

\textbf{指针和字段。} 状态字段要能说清：freq[value] 是频次，heap 元素是 (frequency, value)，heap.top 是淘汰边界。手算时按这个协议跟踪：插入频次 3 和 2 后堆满；频次 1 的元素入堆会被立刻弹出，留下 1 和 2。

\textbf{代码落点。} 关键是比较器按频次建小顶堆。题目通常不要求输出顺序；若要求顺序，最后还要按规则排序结果。关键代码行必须能逐句翻译成“旧状态怎样变成新状态”，否则就只是模板。

\textbf{正确性。} 堆始终保存已处理元素中频次最高的 K 个；任何被弹出的元素频次不高于堆中 K 个候选，不可能进入答案。

\textbf{边界实验。} 先用K 等于不同元素数、频次相同、结果顺序不要求做反例检查。实验时覆盖K 等于不同元素数、频次相同、结果顺序不要求，记录堆顶是否过期、比较器是否让正确候选在堆顶、K 或窗口大小为边界值时是否稳定。

\textbf{复杂度变化。} 完整排序 O(m log m)，m 为不同元素数；堆法 O(n+m log k)，空间 O(m+k)。

\textbf{迁移追问。} 如果题目改成“若要按频次和字典序排序，比较器要同时处理两个键”，先判断原状态是否仍然覆盖新答案集合，再决定是微调边界、改 value 语义，还是换算法。

\noindent\textbf{LeetCode 480 滑动窗口中位数。}

\textbf{题目说明。} LeetCode 480 滑动窗口中位数 的题面入口要先还原成具体任务：窗口移动时既要加入新元素，也要删除滑出的旧元素并查询中位数。

\textbf{样例入口。} 拿 [1, 3, -1]、k=3 手算，中位数是 1；窗口右移时要删除 1、加入 -3，普通堆不能任意删除。

\textbf{暴力基线。} 慢版本每轮扫描全部候选来找最大或最小，再更新候选集合。它的答案选择过程清楚，但动态集合每变化一次就重新找极值，代价太高。放到本题就是：先按“窗口移动时既要加入新元素，也要删除滑出的旧元素并查询中位数”完整试慢版本；样例里已经能看到“规律是双堆维护两半，再用延迟删除清理过期元素，堆顶使用前必须有效”，所以优化前必须先能说清慢版本枚举了哪些候选。

\textbf{暴力过程。} 拿 [1, 3, -1]、k=3 手算，中位数是 1；窗口右移时要删除 1、加入 -3，普通堆不能任意删除。

\textbf{重复工作。} 题面模型：窗口移动时既要加入新元素，也要删除滑出的旧元素并查询中位数。暴力版的慢点要落到本题：慢版本会围绕“窗口移动时既要加入新元素，也要删除滑出的旧元素并查询中位数”反复扫描动态候选集找极值。样例规律是“规律是双堆维护两半，再用延迟删除清理过期元素，堆顶使用前必须有效”，说明每轮真正需要的是堆顶边界，而不是把候选全集重新排序。后面的优化只消掉这类重复工作，不能改变答案集合。

\textbf{优化条件。} 本题的关键观察是：双堆维护左右两半，延迟删除表记录已离窗但尚未到堆顶的元素。这句话必须能翻译成可维护状态；如果题目条件改变后观察不再成立，就不能继续套原来的写法。

\textbf{相邻状态。} 规律是双堆维护两半，再用延迟删除清理过期元素，堆顶使用前必须有效。把这个特例继续拆开看：堆顶必须正好代表下一步最需要处理的候选；若堆里可能有过期对象，读取堆顶前要先清理。堆只解决动态极值，不自动证明被弹出或被丢弃的元素无用。

\textbf{状态复用。} 把每轮全量扫描改成堆顶查询；插入、弹出和过期清理共同维护动态候选集。本题落点是：规律是双堆维护两半，再用延迟删除清理过期元素，堆顶使用前必须有效。手算时只改被新输入影响的那一小部分，而不是回头重算完整候选。

\textbf{指针和字段。} 状态字段要能说清：堆元素字段、堆顶比较规则、候选是否过期、答案读取时机；如果需要懒删除，还要保存删除计数。手算时按这个协议跟踪：拿 [1, 3, -1]、k=3 手算，中位数是 1；窗口右移时要删除 1、加入 -3，普通堆不能任意删除。然后按协议复查：手算时记录堆顶、候选集合和是否有过期元素。每次使用堆顶前都要确认它仍然属于当前问题。

\textbf{代码落点。} 比较器方向先用小样例验证；每次使用堆顶前清理过期项；堆中保存下标时要能回查原数据。关键代码行必须能逐句翻译成“旧状态怎样变成新状态”，否则就只是模板。

\textbf{正确性。} 证明从候选覆盖开始：所有可能成为下一步答案的对象都在堆中，堆顶过期时不会被使用。

\textbf{边界实验。} 先用重复值、偶数窗口、平均值溢出、堆顶过期做反例检查。实验时覆盖重复值、偶数窗口、平均值溢出、堆顶过期，记录堆顶是否过期、比较器是否让正确候选在堆顶、K 或窗口大小为边界值时是否稳定。

\textbf{复杂度变化。} 暴力每轮扫描或排序动态候选，会把线性扫描或排序反复发生；堆把取极值、插入和删除边界控制在对数级。

\textbf{迁移追问。} 如果题目改成“普通 priority\_queue 不能任意删除，所以必须解释懒删除”，先判断原状态是否仍然覆盖新答案集合，再决定是微调边界、改 value 语义，还是换算法。

\noindent\textbf{LeetCode 502 IPO。}

\textbf{题目说明。} LeetCode 502 IPO 的题面入口要先还原成具体任务：当前资本决定可解锁项目集合，目标是在最多 k 次选择中最大化最终资本。

\textbf{样例入口。} 拿 k=2、w=0、profits=[1, 2, 3]、capital=[0, 1, 1] 手算，初始只能做资本 0 的项目，做完利润 1 后资本变 1，才能解锁利润 2 和 3 的项目。

\textbf{暴力基线。} 慢版本每轮扫描全部候选来找最大或最小，再更新候选集合。它的答案选择过程清楚，但动态集合每变化一次就重新找极值，代价太高。放到本题就是：先按“当前资本决定可解锁项目集合，目标是在最多 k 次选择中最大化最终资本”完整试慢版本；样例里已经能看到“规律是按资本排序解锁候选，用利润大顶堆每轮选择当前可做项目里的最大利润”，所以优化前必须先能说清慢版本枚举了哪些候选。

\textbf{暴力过程。} 拿 k=2、w=0、profits=[1, 2, 3]、capital=[0, 1, 1] 手算，初始只能做资本 0 的项目，做完利润 1 后资本变 1，才能解锁利润 2 和 3 的项目。

\textbf{重复工作。} 题面模型：当前资本决定可解锁项目集合，目标是在最多 k 次选择中最大化最终资本。暴力版的慢点要落到本题：慢版本会围绕“当前资本决定可解锁项目集合，目标是在最多 k 次选择中最大化最终资本”反复扫描动态候选集找极值。样例规律是“规律是按资本排序解锁候选，用利润大顶堆每轮选择当前可做项目里的最大利润”，说明每轮真正需要的是堆顶边界，而不是把候选全集重新排序。后面的优化只消掉这类重复工作，不能改变答案集合。

\textbf{优化条件。} 本题的关键观察是：项目按 capital 排序，用指针把可负担项目加入利润大顶堆，每轮选利润最高。这句话必须能翻译成可维护状态；如果题目条件改变后观察不再成立，就不能继续套原来的写法。

\textbf{相邻状态。} 规律是按资本排序解锁候选，用利润大顶堆每轮选择当前可做项目里的最大利润。把这个特例继续拆开看：堆顶必须正好代表下一步最需要处理的候选；若堆里可能有过期对象，读取堆顶前要先清理。堆只解决动态极值，不自动证明被弹出或被丢弃的元素无用。

\textbf{状态复用。} 把每轮全量扫描改成堆顶查询；插入、弹出和过期清理共同维护动态候选集。本题落点是：规律是按资本排序解锁候选，用利润大顶堆每轮选择当前可做项目里的最大利润。手算时只改被新输入影响的那一小部分，而不是回头重算完整候选。

\textbf{指针和字段。} 状态字段要能说清：堆元素字段、堆顶比较规则、候选是否过期、答案读取时机；如果需要懒删除，还要保存删除计数。手算时按这个协议跟踪：拿 k=2、w=0、profits=[1, 2, 3]、capital=[0, 1, 1] 手算，初始只能做资本 0 的项目，做完利润 1 后资本变 1，才能解锁利润 2 和 3 的项目。然后按协议复查：手算时记录堆顶、候选集合和是否有过期元素。每次使用堆顶前都要确认它仍然属于当前问题。

\textbf{代码落点。} 比较器方向先用小样例验证；每次使用堆顶前清理过期项；堆中保存下标时要能回查原数据。关键代码行必须能逐句翻译成“旧状态怎样变成新状态”，否则就只是模板。

\textbf{正确性。} 证明从候选覆盖开始：所有可能成为下一步答案的对象都在堆中，堆顶过期时不会被使用。

\textbf{边界实验。} 先用初始资本为零、没有可做项目、k 大于项目数、利润为零、资本解锁链式增长做反例检查。实验时覆盖初始资本为零、没有可做项目、k 大于项目数、利润为零、资本解锁链式增长，记录堆顶是否过期、比较器是否让正确候选在堆顶、K 或窗口大小为边界值时是否稳定。

\textbf{复杂度变化。} 暴力每轮扫描或排序动态候选，会把线性扫描或排序反复发生；堆把取极值、插入和删除边界控制在对数级。

\textbf{迁移追问。} 如果题目改成“若项目有依赖关系，需要拓扑或约束优化，不能只按资本解锁”，先判断原状态是否仍然覆盖新答案集合，再决定是微调边界、改 value 语义，还是换算法。

\noindent\textbf{LeetCode 621 任务调度器。}

\textbf{题目说明。} LeetCode 621 任务调度器 的题面入口要先还原成具体任务：相同任务之间有冷却间隔，核心受最高频任务骨架限制。

\textbf{样例入口。} 拿 A, A, A, B, B, B、n=2 手算，若只按顺序执行 A 会违反冷却，必须在相同任务之间间隔两个时间单位。

\textbf{暴力基线。} 暴力按时间模拟，每个时刻扫描所有任务，找一个未冷却且剩余次数最多的任务执行，否则空闲。

\textbf{暴力过程。} 任务 A, A, A, B, B, B、n=2 中，相同 A 之间必须间隔两个单位；一种最短安排是 A B idle A B idle A B。

\textbf{重复工作。} 题面模型：相同任务之间有冷却间隔，核心受最高频任务骨架限制。暴力版的慢点要落到本题：浪费在逐时刻扫描所有任务。更重要的是最高频任务决定冷却骨架，其他任务只是填空。后面的优化只消掉这类重复工作，不能改变答案集合。

\textbf{优化条件。} 本题的关键观察是：可用大顶堆模拟，也可用最高频公式计算最短时间。这句话必须能翻译成可维护状态；如果题目条件改变后观察不再成立，就不能继续套原来的写法。

\textbf{相邻状态。} 规律是最高频任务决定骨架长度，也可用堆按剩余次数模拟，每轮安排最多 n+1 个不同任务。把这个特例继续拆开看：堆顶必须正好代表下一步最需要处理的候选；若堆里可能有过期对象，读取堆顶前要先清理。堆只解决动态极值，不自动证明被弹出或被丢弃的元素无用。

\textbf{状态复用。} 公式法统计最大频次 maxFreq 和拥有最大频次的任务数 maxCount。骨架长度为 (maxFreq-1)*(n+1)+maxCount，答案取它和任务总数的较大值。手算时只改被新输入影响的那一小部分，而不是回头重算完整候选。

\textbf{指针和字段。} 状态字段要能说清：maxFreq 是最高频任务次数，maxCount 是最高频任务种类数，n 是冷却间隔，tasks.size() 是没有空闲时的下界。手算时按这个协议跟踪：A 和 B 都出现 3 次，maxFreq=3、maxCount=2，骨架为 (3-1)*3+2=8。

\textbf{代码落点。} 关键是最后 return max(tasks.size(), skeleton)。若 n=0，任务总数就是答案。关键代码行必须能逐句翻译成“旧状态怎样变成新状态”，否则就只是模板。

\textbf{正确性。} 最高频任务之间至少有 maxFreq-1 个间隔块，每块长度 n+1；多个最高频任务占据最后一列。其他任务只能填充空位，不能突破任务总数下界。

\textbf{边界实验。} 先用冷却为零、多个最高频、空闲被其他任务填满做反例检查。实验时覆盖冷却为零、多个最高频、空闲被其他任务填满，记录堆顶是否过期、比较器是否让正确候选在堆顶、K 或窗口大小为边界值时是否稳定。

\textbf{复杂度变化。} 堆模拟 O(T log C)，公式法统计频次 O(T+C)，空间 O(C)，C 为任务种类数。

\textbf{迁移追问。} 如果题目改成“若任务有不同释放时间或执行时长，就变成更一般的调度问题”，先判断原状态是否仍然覆盖新答案集合，再决定是微调边界、改 value 语义，还是换算法。

\noindent\textbf{LeetCode 632 最小区间覆盖 K 个列表。}

\textbf{题目说明。} LeetCode 632 最小区间覆盖 K 个列表 的题面入口要先还原成具体任务：每个列表至少取一个元素时，当前区间由最小当前值和最大当前值决定。

\textbf{样例入口。} 拿三组 [4, 10]、[0, 9]、[5, 18] 手算，当前头是 4、0、5，区间 [0, 5] 覆盖三组；弹出最小头 0 后推进该组。

\textbf{暴力基线。} 慢版本每轮扫描全部候选来找最大或最小，再更新候选集合。它的答案选择过程清楚，但动态集合每变化一次就重新找极值，代价太高。放到本题就是：先按“每个列表至少取一个元素时，当前区间由最小当前值和最大当前值决定”完整试慢版本；样例里已经能看到“规律是堆保存每组当前元素，同时维护当前最大值，最小头决定下一步推进哪一组”，所以优化前必须先能说清慢版本枚举了哪些候选。

\textbf{暴力过程。} 拿三组 [4, 10]、[0, 9]、[5, 18] 手算，当前头是 4、0、5，区间 [0, 5] 覆盖三组；弹出最小头 0 后推进该组。

\textbf{重复工作。} 题面模型：每个列表至少取一个元素时，当前区间由最小当前值和最大当前值决定。暴力版的慢点要落到本题：慢版本会围绕“每个列表至少取一个元素时，当前区间由最小当前值和最大当前值决定”反复扫描动态候选集找极值。样例规律是“规律是堆保存每组当前元素，同时维护当前最大值，最小头决定下一步推进哪一组”，说明每轮真正需要的是堆顶边界，而不是把候选全集重新排序。后面的优化只消掉这类重复工作，不能改变答案集合。

\textbf{优化条件。} 本题的关键观察是：小顶堆保存各列表当前元素，同时维护当前最大值；弹出最小所在列表并推进。这句话必须能翻译成可维护状态；如果题目条件改变后观察不再成立，就不能继续套原来的写法。

\textbf{相邻状态。} 规律是堆保存每组当前元素，同时维护当前最大值，最小头决定下一步推进哪一组。把这个特例继续拆开看：堆顶必须正好代表下一步最需要处理的候选；若堆里可能有过期对象，读取堆顶前要先清理。堆只解决动态极值，不自动证明被弹出或被丢弃的元素无用。

\textbf{状态复用。} 把每轮全量扫描改成堆顶查询；插入、弹出和过期清理共同维护动态候选集。本题落点是：规律是堆保存每组当前元素，同时维护当前最大值，最小头决定下一步推进哪一组。手算时只改被新输入影响的那一小部分，而不是回头重算完整候选。

\textbf{指针和字段。} 状态字段要能说清：堆元素字段、堆顶比较规则、候选是否过期、答案读取时机；如果需要懒删除，还要保存删除计数。手算时按这个协议跟踪：拿三组 [4, 10]、[0, 9]、[5, 18] 手算，当前头是 4、0、5，区间 [0, 5] 覆盖三组；弹出最小头 0 后推进该组。然后按协议复查：手算时记录堆顶、候选集合和是否有过期元素。每次使用堆顶前都要确认它仍然属于当前问题。

\textbf{代码落点。} 比较器方向先用小样例验证；每次使用堆顶前清理过期项；堆中保存下标时要能回查原数据。关键代码行必须能逐句翻译成“旧状态怎样变成新状态”，否则就只是模板。

\textbf{正确性。} 证明从候选覆盖开始：所有可能成为下一步答案的对象都在堆中，堆顶过期时不会被使用。

\textbf{边界实验。} 先用某个列表为空、重复值、区间长度相同的 tie-breaker、列表耗尽做反例检查。实验时覆盖某个列表为空、重复值、区间长度相同的 tie-breaker、列表耗尽，记录堆顶是否过期、比较器是否让正确候选在堆顶、K 或窗口大小为边界值时是否稳定。

\textbf{复杂度变化。} 暴力每轮扫描或排序动态候选，会把线性扫描或排序反复发生；堆把取极值、插入和删除边界控制在对数级。

\textbf{迁移追问。} 如果题目改成“这是多路归并和覆盖窗口的结合”，先判断原状态是否仍然覆盖新答案集合，再决定是微调边界、改 value 语义，还是换算法。

\noindent\textbf{LeetCode 703 数据流中的第 K 大元素。}

\textbf{题目说明。} LeetCode 703 数据流中的第 K 大元素 的题面入口要先还原成具体任务：数据不断到达，只关心当前第 K 大而非完整排序。

\textbf{样例入口。} 拿 k=3，初始 4, 5, 8, 2 手算，小顶堆保存 4, 5, 8，堆顶 4 是第 3 大；add 3 不进入堆，add 5 会替换 4。

\textbf{暴力基线。} 慢版本每轮扫描全部候选来找最大或最小，再更新候选集合。它的答案选择过程清楚，但动态集合每变化一次就重新找极值，代价太高。放到本题就是：先按“数据不断到达，只关心当前第 K 大而非完整排序”完整试慢版本；样例里已经能看到“规律是数据流只保留前 k 大，堆顶就是当前第 k 大”，所以优化前必须先能说清慢版本枚举了哪些候选。

\textbf{暴力过程。} 拿 k=3，初始 4, 5, 8, 2 手算，小顶堆保存 4, 5, 8，堆顶 4 是第 3 大；add 3 不进入堆，add 5 会替换 4。

\textbf{重复工作。} 题面模型：数据不断到达，只关心当前第 K 大而非完整排序。暴力版的慢点要落到本题：慢版本会围绕“数据不断到达，只关心当前第 K 大而非完整排序”反复扫描动态候选集找极值。样例规律是“规律是数据流只保留前 k 大，堆顶就是当前第 k 大”，说明每轮真正需要的是堆顶边界，而不是把候选全集重新排序。后面的优化只消掉这类重复工作，不能改变答案集合。

\textbf{优化条件。} 本题的关键观察是：维护大小为 K 的小顶堆，堆顶就是当前第 K 大元素。这句话必须能翻译成可维护状态；如果题目条件改变后观察不再成立，就不能继续套原来的写法。

\textbf{相邻状态。} 规律是数据流只保留前 k 大，堆顶就是当前第 k 大。把这个特例继续拆开看：堆顶必须正好代表下一步最需要处理的候选；若堆里可能有过期对象，读取堆顶前要先清理。堆只解决动态极值，不自动证明被弹出或被丢弃的元素无用。

\textbf{状态复用。} 把每轮全量扫描改成堆顶查询；插入、弹出和过期清理共同维护动态候选集。本题落点是：规律是数据流只保留前 k 大，堆顶就是当前第 k 大。手算时只改被新输入影响的那一小部分，而不是回头重算完整候选。

\textbf{指针和字段。} 状态字段要能说清：堆元素字段、堆顶比较规则、候选是否过期、答案读取时机；如果需要懒删除，还要保存删除计数。手算时按这个协议跟踪：拿 k=3，初始 4, 5, 8, 2 手算，小顶堆保存 4, 5, 8，堆顶 4 是第 3 大；add 3 不进入堆，add 5 会替换 4。然后按协议复查：手算时记录堆顶、候选集合和是否有过期元素。每次使用堆顶前都要确认它仍然属于当前问题。

\textbf{代码落点。} 比较器方向先用小样例验证；每次使用堆顶前清理过期项；堆中保存下标时要能回查原数据。关键代码行必须能逐句翻译成“旧状态怎样变成新状态”，否则就只是模板。

\textbf{正确性。} 证明从候选覆盖开始：所有可能成为下一步答案的对象都在堆中，堆顶过期时不会被使用。

\textbf{边界实验。} 先用初始元素少于 K、重复值、K 为一、连续 add做反例检查。实验时覆盖初始元素少于 K、重复值、K 为一、连续 add，记录堆顶是否过期、比较器是否让正确候选在堆顶、K 或窗口大小为边界值时是否稳定。

\textbf{复杂度变化。} 暴力每轮扫描或排序动态候选，会把线性扫描或排序反复发生；堆把取极值、插入和删除边界控制在对数级。

\textbf{迁移追问。} 如果题目改成“这题是 Top K 的在线版本，和一次性排序的适用场景不同”，先判断原状态是否仍然覆盖新答案集合，再决定是微调边界、改 value 语义，还是换算法。

\noindent\textbf{LeetCode 767 重构字符串。}

\textbf{题目说明。} LeetCode 767 重构字符串 的题面入口要先还原成具体任务：相邻不能相同，剩余次数最多的字符最难安排。

\textbf{样例入口。} 拿 s=aab 手算，a 剩 2 次、b 剩 1 次，每轮取两个最高频字符放入答案，可得到 ab 后再放回剩余 a，最后是 aba。拿 aaab 时 a 出现 3 次超过 (4+1)/2，不可能隔开。

\textbf{暴力基线。} 慢版本每轮扫描全部候选来找最大或最小，再更新候选集合。它的答案选择过程清楚，但动态集合每变化一次就重新找极值，代价太高。放到本题就是：先按“相邻不能相同，剩余次数最多的字符最难安排”完整试慢版本；样例里已经能看到“规律是最大频次先判不可能，构造时堆里每轮取两个不同最高频字符”，所以优化前必须先能说清慢版本枚举了哪些候选。

\textbf{暴力过程。} 拿 s=aab 手算，a 剩 2 次、b 剩 1 次，每轮取两个最高频字符放入答案，可得到 ab 后再放回剩余 a，最后是 aba。拿 aaab 时 a 出现 3 次超过 (4+1)/2，不可能隔开。

\textbf{重复工作。} 题面模型：相邻不能相同，剩余次数最多的字符最难安排。暴力版的慢点要落到本题：慢版本会围绕“相邻不能相同，剩余次数最多的字符最难安排”反复扫描动态候选集找极值。样例规律是“规律是最大频次先判不可能，构造时堆里每轮取两个不同最高频字符”，说明每轮真正需要的是堆顶边界，而不是把候选全集重新排序。后面的优化只消掉这类重复工作，不能改变答案集合。

\textbf{优化条件。} 本题的关键观察是：大顶堆每次取两个最高频不同字符放入答案，剩余频次再放回；先检查最大频次是否超过 (n+1)/2。这句话必须能翻译成可维护状态；如果题目条件改变后观察不再成立，就不能继续套原来的写法。

\textbf{相邻状态。} 规律是最大频次先判不可能，构造时堆里每轮取两个不同最高频字符。把这个特例继续拆开看：堆顶必须正好代表下一步最需要处理的候选；若堆里可能有过期对象，读取堆顶前要先清理。堆只解决动态极值，不自动证明被弹出或被丢弃的元素无用。

\textbf{状态复用。} 把每轮全量扫描改成堆顶查询；插入、弹出和过期清理共同维护动态候选集。本题落点是：规律是最大频次先判不可能，构造时堆里每轮取两个不同最高频字符。手算时只改被新输入影响的那一小部分，而不是回头重算完整候选。

\textbf{指针和字段。} 状态字段要能说清：堆元素字段、堆顶比较规则、候选是否过期、答案读取时机；如果需要懒删除，还要保存删除计数。手算时按这个协议跟踪：拿 s=aab 手算，a 剩 2 次、b 剩 1 次，每轮取两个最高频字符放入答案，可得到 ab 后再放回剩余 a，最后是 aba。拿 aaab 时 a 出现 3 次超过 (4+1)/2，不可能隔开。然后按协议复查：手算时记录堆顶、候选集合和是否有过期元素。每次使用堆顶前都要确认它仍然属于当前问题。

\textbf{代码落点。} 比较器方向先用小样例验证；每次使用堆顶前清理过期项；堆中保存下标时要能回查原数据。关键代码行必须能逐句翻译成“旧状态怎样变成新状态”，否则就只是模板。

\textbf{正确性。} 证明从候选覆盖开始：所有可能成为下一步答案的对象都在堆中，堆顶过期时不会被使用。

\textbf{边界实验。} 先用单字符、最高频过半、多个字符同频、最后只剩一个字符、结果不唯一做反例检查。实验时覆盖单字符、最高频过半、多个字符同频、最后只剩一个字符、结果不唯一，记录堆顶是否过期、比较器是否让正确候选在堆顶、K 或窗口大小为边界值时是否稳定。

\textbf{复杂度变化。} 暴力每轮扫描或排序动态候选，会把线性扫描或排序反复发生；堆把取极值、插入和删除边界控制在对数级。

\textbf{迁移追问。} 如果题目改成“若要求相同字符至少间隔 k，堆还要配合冷却队列”，先判断原状态是否仍然覆盖新答案集合，再决定是微调边界、改 value 语义，还是换算法。

\noindent\textbf{LeetCode 857 雇佣 K 名工人的最低成本。}

\textbf{题目说明。} LeetCode 857 雇佣 K 名工人的最低成本 的题面入口要先还原成具体任务：工资比例由当前队伍中最高工资质量比决定，总成本等于比例乘质量和。

\textbf{样例入口。} 拿 quality=[10, 20, 5]、wage=[70, 50, 30] 手算，按工资质量比选组时，当前最大比率决定组内所有人的单位工资；再用堆保留质量和最小的 k 人。

\textbf{暴力基线。} 慢版本每轮扫描全部候选来找最大或最小，再更新候选集合。它的答案选择过程清楚，但动态集合每变化一次就重新找极值，代价太高。放到本题就是：先按“工资比例由当前队伍中最高工资质量比决定，总成本等于比例乘质量和”完整试慢版本；样例里已经能看到“规律是排序枚举最高比率，堆优化质量总和”，所以优化前必须先能说清慢版本枚举了哪些候选。

\textbf{暴力过程。} 拿 quality=[10, 20, 5]、wage=[70, 50, 30] 手算，按工资质量比选组时，当前最大比率决定组内所有人的单位工资；再用堆保留质量和最小的 k 人。

\textbf{重复工作。} 题面模型：工资比例由当前队伍中最高工资质量比决定，总成本等于比例乘质量和。暴力版的慢点要落到本题：慢版本会围绕“工资比例由当前队伍中最高工资质量比决定，总成本等于比例乘质量和”反复扫描动态候选集找极值。样例规律是“规律是排序枚举最高比率，堆优化质量总和”，说明每轮真正需要的是堆顶边界，而不是把候选全集重新排序。后面的优化只消掉这类重复工作，不能改变答案集合。

\textbf{优化条件。} 本题的关键观察是：按比例升序扫描，用大顶堆维护当前 K 个工人的最小质量和。这句话必须能翻译成可维护状态；如果题目条件改变后观察不再成立，就不能继续套原来的写法。

\textbf{相邻状态。} 规律是排序枚举最高比率，堆优化质量总和。把这个特例继续拆开看：堆顶必须正好代表下一步最需要处理的候选；若堆里可能有过期对象，读取堆顶前要先清理。堆只解决动态极值，不自动证明被弹出或被丢弃的元素无用。

\textbf{状态复用。} 把每轮全量扫描改成堆顶查询；插入、弹出和过期清理共同维护动态候选集。本题落点是：规律是排序枚举最高比率，堆优化质量总和。手算时只改被新输入影响的那一小部分，而不是回头重算完整候选。

\textbf{指针和字段。} 状态字段要能说清：堆元素字段、堆顶比较规则、候选是否过期、答案读取时机；如果需要懒删除，还要保存删除计数。手算时按这个协议跟踪：拿 quality=[10, 20, 5]、wage=[70, 50, 30] 手算，按工资质量比选组时，当前最大比率决定组内所有人的单位工资；再用堆保留质量和最小的 k 人。然后按协议复查：手算时记录堆顶、候选集合和是否有过期元素。每次使用堆顶前都要确认它仍然属于当前问题。

\textbf{代码落点。} 比较器方向先用小样例验证；每次使用堆顶前清理过期项；堆中保存下标时要能回查原数据。关键代码行必须能逐句翻译成“旧状态怎样变成新状态”，否则就只是模板。

\textbf{正确性。} 证明从候选覆盖开始：所有可能成为下一步答案的对象都在堆中，堆顶过期时不会被使用。

\textbf{边界实验。} 先用K 为一、比例相同、浮点精度、质量和更新做反例检查。实验时覆盖K 为一、比例相同、浮点精度、质量和更新，记录堆顶是否过期、比较器是否让正确候选在堆顶、K 或窗口大小为边界值时是否稳定。

\textbf{复杂度变化。} 暴力每轮扫描或排序动态候选，会把线性扫描或排序反复发生；堆把取极值、插入和删除边界控制在对数级。

\textbf{迁移追问。} 如果题目改成“这是排序确定价格比例，堆维护可替换候选的组合题”，先判断原状态是否仍然覆盖新答案集合，再决定是微调边界、改 value 语义，还是换算法。

\noindent\textbf{LeetCode 871 最低加油次数。}

\textbf{题目说明。} LeetCode 871 最低加油次数 的题面入口要先还原成具体任务：行驶过程中，经过的加油站形成可反悔的燃料候选。

\textbf{样例入口。} 拿 target=100、startFuel=10 手算，车先到 10 号站，把 60 升油加入最大堆；之后油量不足时，从已经路过的站里取当前最大油量。

\textbf{暴力基线。} 慢版本每轮扫描全部候选来找最大或最小，再更新候选集合。它的答案选择过程清楚，但动态集合每变化一次就重新找极值，代价太高。放到本题就是：先按“行驶过程中，经过的加油站形成可反悔的燃料候选”完整试慢版本；样例里已经能看到“规律是能到的站先入堆，需要加油时再反悔选择最大燃料，保证次数最少”，所以优化前必须先能说清慢版本枚举了哪些候选。

\textbf{暴力过程。} 拿 target=100、startFuel=10 手算，车先到 10 号站，把 60 升油加入最大堆；之后油量不足时，从已经路过的站里取当前最大油量。

\textbf{重复工作。} 题面模型：行驶过程中，经过的加油站形成可反悔的燃料候选。暴力版的慢点要落到本题：慢版本会围绕“行驶过程中，经过的加油站形成可反悔的燃料候选”反复扫描动态候选集找极值。样例规律是“规律是能到的站先入堆，需要加油时再反悔选择最大燃料，保证次数最少”，说明每轮真正需要的是堆顶边界，而不是把候选全集重新排序。后面的优化只消掉这类重复工作，不能改变答案集合。

\textbf{优化条件。} 本题的关键观察是：按位置扫描，油量不足到达下一点时，从已路过站点中取最大燃料补充。这句话必须能翻译成可维护状态；如果题目条件改变后观察不再成立，就不能继续套原来的写法。

\textbf{相邻状态。} 规律是能到的站先入堆，需要加油时再反悔选择最大燃料，保证次数最少。把这个特例继续拆开看：堆顶必须正好代表下一步最需要处理的候选；若堆里可能有过期对象，读取堆顶前要先清理。堆只解决动态极值，不自动证明被弹出或被丢弃的元素无用。

\textbf{状态复用。} 把每轮全量扫描改成堆顶查询；插入、弹出和过期清理共同维护动态候选集。本题落点是：规律是能到的站先入堆，需要加油时再反悔选择最大燃料，保证次数最少。手算时只改被新输入影响的那一小部分，而不是回头重算完整候选。

\textbf{指针和字段。} 状态字段要能说清：堆元素字段、堆顶比较规则、候选是否过期、答案读取时机；如果需要懒删除，还要保存删除计数。手算时按这个协议跟踪：拿 target=100、startFuel=10 手算，车先到 10 号站，把 60 升油加入最大堆；之后油量不足时，从已经路过的站里取当前最大油量。然后按协议复查：手算时记录堆顶、候选集合和是否有过期元素。每次使用堆顶前都要确认它仍然属于当前问题。

\textbf{代码落点。} 比较器方向先用小样例验证；每次使用堆顶前清理过期项；堆中保存下标时要能回查原数据。关键代码行必须能逐句翻译成“旧状态怎样变成新状态”，否则就只是模板。

\textbf{正确性。} 证明从候选覆盖开始：所有可能成为下一步答案的对象都在堆中，堆顶过期时不会被使用。

\textbf{边界实验。} 先用起点油足够、无法到达第一个站、多个站同位置、目标作为终点事件做反例检查。实验时覆盖起点油足够、无法到达第一个站、多个站同位置、目标作为终点事件，记录堆顶是否过期、比较器是否让正确候选在堆顶、K 或窗口大小为边界值时是否稳定。

\textbf{复杂度变化。} 暴力每轮扫描或排序动态候选，会把线性扫描或排序反复发生；堆把取极值、插入和删除边界控制在对数级。

\textbf{迁移追问。} 如果题目改成“这是反悔贪心，堆里保存的是已经经过但尚未使用的资源”，先判断原状态是否仍然覆盖新答案集合，再决定是微调边界、改 value 语义，还是换算法。

\noindent\textbf{LeetCode 973 最接近原点的 K 个点。}

\textbf{题目说明。} LeetCode 973 最接近原点的 K 个点 的题面入口要先还原成具体任务：距离只需比较平方，避免开方误差和成本。

\textbf{样例入口。} 拿点 (1, 3) 和 (-2, 2) 手算，距离平方分别是 10 和 8，所以 k=1 时选 (-2, 2)。

\textbf{暴力基线。} 慢版本每轮扫描全部候选来找最大或最小，再更新候选集合。它的答案选择过程清楚，但动态集合每变化一次就重新找极值，代价太高。放到本题就是：先按“距离只需比较平方，避免开方误差和成本”完整试慢版本；样例里已经能看到“规律是比较距离平方避免开方，用大小为 k 的大顶堆保留当前最近点”，所以优化前必须先能说清慢版本枚举了哪些候选。

\textbf{暴力过程。} 拿点 (1, 3) 和 (-2, 2) 手算，距离平方分别是 10 和 8，所以 k=1 时选 (-2, 2)。

\textbf{重复工作。} 题面模型：距离只需比较平方，避免开方误差和成本。暴力版的慢点要落到本题：慢版本会围绕“距离只需比较平方，避免开方误差和成本”反复扫描动态候选集找极值。样例规律是“规律是比较距离平方避免开方，用大小为 k 的大顶堆保留当前最近点”，说明每轮真正需要的是堆顶边界，而不是把候选全集重新排序。后面的优化只消掉这类重复工作，不能改变答案集合。

\textbf{优化条件。} 本题的关键观察是：大小为 K 的大顶堆保存当前最近 K 个点中最远者。这句话必须能翻译成可维护状态；如果题目条件改变后观察不再成立，就不能继续套原来的写法。

\textbf{相邻状态。} 规律是比较距离平方避免开方，用大小为 k 的大顶堆保留当前最近点。把这个特例继续拆开看：堆顶必须正好代表下一步最需要处理的候选；若堆里可能有过期对象，读取堆顶前要先清理。堆只解决动态极值，不自动证明被弹出或被丢弃的元素无用。

\textbf{状态复用。} 把每轮全量扫描改成堆顶查询；插入、弹出和过期清理共同维护动态候选集。本题落点是：规律是比较距离平方避免开方，用大小为 k 的大顶堆保留当前最近点。手算时只改被新输入影响的那一小部分，而不是回头重算完整候选。

\textbf{指针和字段。} 状态字段要能说清：堆元素字段、堆顶比较规则、候选是否过期、答案读取时机；如果需要懒删除，还要保存删除计数。手算时按这个协议跟踪：拿点 (1, 3) 和 (-2, 2) 手算，距离平方分别是 10 和 8，所以 k=1 时选 (-2, 2)。然后按协议复查：手算时记录堆顶、候选集合和是否有过期元素。每次使用堆顶前都要确认它仍然属于当前问题。

\textbf{代码落点。} 比较器方向先用小样例验证；每次使用堆顶前清理过期项；堆中保存下标时要能回查原数据。关键代码行必须能逐句翻译成“旧状态怎样变成新状态”，否则就只是模板。

\textbf{正确性。} 证明从候选覆盖开始：所有可能成为下一步答案的对象都在堆中，堆顶过期时不会被使用。

\textbf{边界实验。} 先用K 为一、K 为全部、距离相同、坐标平方溢出做反例检查。实验时覆盖K 为一、K 为全部、距离相同、坐标平方溢出，记录堆顶是否过期、比较器是否让正确候选在堆顶、K 或窗口大小为边界值时是否稳定。

\textbf{复杂度变化。} 暴力每轮扫描或排序动态候选，会把线性扫描或排序反复发生；堆把取极值、插入和删除边界控制在对数级。

\textbf{迁移追问。} 如果题目改成“快速选择可平均线性，但堆更适合数据流”，先判断原状态是否仍然覆盖新答案集合，再决定是微调边界、改 value 语义，还是换算法。

\noindent\textbf{LeetCode 1046 最后一块石头的重量。}

\textbf{题目说明。} LeetCode 1046 最后一块石头的重量 的题面入口要先还原成具体任务：每次都要取当前最重的两块石头碰撞。

\textbf{样例入口。} 拿 [2, 7, 4, 1, 8, 1] 手算，每次取两块最大石头 8 和 7，撞剩 1 再放回。

\textbf{暴力基线。} 慢版本每轮扫描全部候选来找最大或最小，再更新候选集合。它的答案选择过程清楚，但动态集合每变化一次就重新找极值，代价太高。放到本题就是：先按“每次都要取当前最重的两块石头碰撞”完整试慢版本；样例里已经能看到“规律是动态最大值查询，用大顶堆保存当前石头重量，直到剩 0 或 1 块”，所以优化前必须先能说清慢版本枚举了哪些候选。

\textbf{暴力过程。} 拿 [2, 7, 4, 1, 8, 1] 手算，每次取两块最大石头 8 和 7，撞剩 1 再放回。

\textbf{重复工作。} 题面模型：每次都要取当前最重的两块石头碰撞。暴力版的慢点要落到本题：慢版本会围绕“每次都要取当前最重的两块石头碰撞”反复扫描动态候选集找极值。样例规律是“规律是动态最大值查询，用大顶堆保存当前石头重量，直到剩 0 或 1 块”，说明每轮真正需要的是堆顶边界，而不是把候选全集重新排序。后面的优化只消掉这类重复工作，不能改变答案集合。

\textbf{优化条件。} 本题的关键观察是：大顶堆保存所有重量，弹出两个最大值，若不同则把差值放回。这句话必须能翻译成可维护状态；如果题目条件改变后观察不再成立，就不能继续套原来的写法。

\textbf{相邻状态。} 规律是动态最大值查询，用大顶堆保存当前石头重量，直到剩 0 或 1 块。把这个特例继续拆开看：堆顶必须正好代表下一步最需要处理的候选；若堆里可能有过期对象，读取堆顶前要先清理。堆只解决动态极值，不自动证明被弹出或被丢弃的元素无用。

\textbf{状态复用。} 把每轮全量扫描改成堆顶查询；插入、弹出和过期清理共同维护动态候选集。本题落点是：规律是动态最大值查询，用大顶堆保存当前石头重量，直到剩 0 或 1 块。手算时只改被新输入影响的那一小部分，而不是回头重算完整候选。

\textbf{指针和字段。} 状态字段要能说清：堆元素字段、堆顶比较规则、候选是否过期、答案读取时机；如果需要懒删除，还要保存删除计数。手算时按这个协议跟踪：拿 [2, 7, 4, 1, 8, 1] 手算，每次取两块最大石头 8 和 7，撞剩 1 再放回。然后按协议复查：手算时记录堆顶、候选集合和是否有过期元素。每次使用堆顶前都要确认它仍然属于当前问题。

\textbf{代码落点。} 比较器方向先用小样例验证；每次使用堆顶前清理过期项；堆中保存下标时要能回查原数据。关键代码行必须能逐句翻译成“旧状态怎样变成新状态”，否则就只是模板。

\textbf{正确性。} 证明从候选覆盖开始：所有可能成为下一步答案的对象都在堆中，堆顶过期时不会被使用。

\textbf{边界实验。} 先用没有石头、一块石头、两块相等、多次产生相同重量做反例检查。实验时覆盖没有石头、一块石头、两块相等、多次产生相同重量，记录堆顶是否过期、比较器是否让正确候选在堆顶、K 或窗口大小为边界值时是否稳定。

\textbf{复杂度变化。} 暴力每轮扫描或排序动态候选，会把线性扫描或排序反复发生；堆把取极值、插入和删除边界控制在对数级。

\textbf{迁移追问。} 如果题目改成“最后一块石头 II 则是分组差值最小化，模型转为 0-1 背包”，先判断原状态是否仍然覆盖新答案集合，再决定是微调边界、改 value 语义，还是换算法。

\noindent\textbf{LeetCode 1353 最多可以参加的会议数目。}

\textbf{题目说明。} LeetCode 1353 最多可以参加的会议数目 的题面入口要先还原成具体任务：每天只能参加一个正在进行的会议，应优先选择最早结束的会议。

\textbf{样例入口。} 拿 events=[[1, 2], [2, 3], [3, 4]] 手算，第 1 天把开始的会议入堆并参加最早结束的；第 2 天再清理过期并参加一个。

\textbf{暴力基线。} 慢版本每轮扫描全部候选来找最大或最小，再更新候选集合。它的答案选择过程清楚，但动态集合每变化一次就重新找极值，代价太高。放到本题就是：先按“每天只能参加一个正在进行的会议，应优先选择最早结束的会议”完整试慢版本；样例里已经能看到“规律是按天扫描，用小顶堆保存当前可参加会议的结束日，每天选最早结束者”，所以优化前必须先能说清慢版本枚举了哪些候选。

\textbf{暴力过程。} 拿 events=[[1, 2], [2, 3], [3, 4]] 手算，第 1 天把开始的会议入堆并参加最早结束的；第 2 天再清理过期并参加一个。

\textbf{重复工作。} 题面模型：每天只能参加一个正在进行的会议，应优先选择最早结束的会议。暴力版的慢点要落到本题：慢版本会围绕“每天只能参加一个正在进行的会议，应优先选择最早结束的会议”反复扫描动态候选集找极值。样例规律是“规律是按天扫描，用小顶堆保存当前可参加会议的结束日，每天选最早结束者”，说明每轮真正需要的是堆顶边界，而不是把候选全集重新排序。后面的优化只消掉这类重复工作，不能改变答案集合。

\textbf{优化条件。} 本题的关键观察是：按开始日排序扫描日期，把当天可参加会议的结束日入小顶堆，弹出最早结束者。这句话必须能翻译成可维护状态；如果题目条件改变后观察不再成立，就不能继续套原来的写法。

\textbf{相邻状态。} 规律是按天扫描，用小顶堆保存当前可参加会议的结束日，每天选最早结束者。把这个特例继续拆开看：堆顶必须正好代表下一步最需要处理的候选；若堆里可能有过期对象，读取堆顶前要先清理。堆只解决动态极值，不自动证明被弹出或被丢弃的元素无用。

\textbf{状态复用。} 把每轮全量扫描改成堆顶查询；插入、弹出和过期清理共同维护动态候选集。本题落点是：规律是按天扫描，用小顶堆保存当前可参加会议的结束日，每天选最早结束者。手算时只改被新输入影响的那一小部分，而不是回头重算完整候选。

\textbf{指针和字段。} 状态字段要能说清：堆元素字段、堆顶比较规则、候选是否过期、答案读取时机；如果需要懒删除，还要保存删除计数。手算时按这个协议跟踪：拿 events=[[1, 2], [2, 3], [3, 4]] 手算，第 1 天把开始的会议入堆并参加最早结束的；第 2 天再清理过期并参加一个。然后按协议复查：手算时记录堆顶、候选集合和是否有过期元素。每次使用堆顶前都要确认它仍然属于当前问题。

\textbf{代码落点。} 比较器方向先用小样例验证；每次使用堆顶前清理过期项；堆中保存下标时要能回查原数据。关键代码行必须能逐句翻译成“旧状态怎样变成新状态”，否则就只是模板。

\textbf{正确性。} 证明从候选覆盖开始：所有可能成为下一步答案的对象都在堆中，堆顶过期时不会被使用。

\textbf{边界实验。} 先用同一天多个会议、过期会议清理、空闲日期跳跃、结束日相同做反例检查。实验时覆盖同一天多个会议、过期会议清理、空闲日期跳跃、结束日相同，记录堆顶是否过期、比较器是否让正确候选在堆顶、K 或窗口大小为边界值时是否稳定。

\textbf{复杂度变化。} 暴力每轮扫描或排序动态候选，会把线性扫描或排序反复发生；堆把取极值、插入和删除边界控制在对数级。

\textbf{迁移追问。} 如果题目改成“这是扫描线和堆结合的调度题”，先判断原状态是否仍然覆盖新答案集合，再决定是微调边界、改 value 语义，还是换算法。

\noindent\textbf{LeetCode 1675 数组的最小偏移量。}

\textbf{题目说明。} LeetCode 1675 数组的最小偏移量 的题面入口要先还原成具体任务：奇数可以先乘二，偶数可以不断除二，目标是缩小当前最大最小差。

\textbf{样例入口。} 拿 [1, 2, 3, 4] 手算，奇数可先翻倍成偶数，之后只能不断把当前最大偶数减半；每次记录最大值和当前最小值差。

\textbf{暴力基线。} 慢版本每轮扫描全部候选来找最大或最小，再更新候选集合。它的答案选择过程清楚，但动态集合每变化一次就重新找极值，代价太高。放到本题就是：先按“奇数可以先乘二，偶数可以不断除二，目标是缩小当前最大最小差”完整试慢版本；样例里已经能看到“规律是把所有数变到可下降的最高状态，再用大顶堆逐步降低最大值”，所以优化前必须先能说清慢版本枚举了哪些候选。

\textbf{暴力过程。} 拿 [1, 2, 3, 4] 手算，奇数可先翻倍成偶数，之后只能不断把当前最大偶数减半；每次记录最大值和当前最小值差。

\textbf{重复工作。} 题面模型：奇数可以先乘二，偶数可以不断除二，目标是缩小当前最大最小差。暴力版的慢点要落到本题：慢版本会围绕“奇数可以先乘二，偶数可以不断除二，目标是缩小当前最大最小差”反复扫描动态候选集找极值。样例规律是“规律是把所有数变到可下降的最高状态，再用大顶堆逐步降低最大值”，说明每轮真正需要的是堆顶边界，而不是把候选全集重新排序。后面的优化只消掉这类重复工作，不能改变答案集合。

\textbf{优化条件。} 本题的关键观察是：把所有数规范到可降低的偶数形态，用大顶堆维护当前最大值，同时记录当前最小值。这句话必须能翻译成可维护状态；如果题目条件改变后观察不再成立，就不能继续套原来的写法。

\textbf{相邻状态。} 规律是把所有数变到可下降的最高状态，再用大顶堆逐步降低最大值。把这个特例继续拆开看：堆顶必须正好代表下一步最需要处理的候选；若堆里可能有过期对象，读取堆顶前要先清理。堆只解决动态极值，不自动证明被弹出或被丢弃的元素无用。

\textbf{状态复用。} 把每轮全量扫描改成堆顶查询；插入、弹出和过期清理共同维护动态候选集。本题落点是：规律是把所有数变到可下降的最高状态，再用大顶堆逐步降低最大值。手算时只改被新输入影响的那一小部分，而不是回头重算完整候选。

\textbf{指针和字段。} 状态字段要能说清：堆元素字段、堆顶比较规则、候选是否过期、答案读取时机；如果需要懒删除，还要保存删除计数。手算时按这个协议跟踪：拿 [1, 2, 3, 4] 手算，奇数可先翻倍成偶数，之后只能不断把当前最大偶数减半；每次记录最大值和当前最小值差。然后按协议复查：手算时记录堆顶、候选集合和是否有过期元素。每次使用堆顶前都要确认它仍然属于当前问题。

\textbf{代码落点。} 比较器方向先用小样例验证；每次使用堆顶前清理过期项；堆中保存下标时要能回查原数据。关键代码行必须能逐句翻译成“旧状态怎样变成新状态”，否则就只是模板。

\textbf{正确性。} 证明从候选覆盖开始：所有可能成为下一步答案的对象都在堆中，堆顶过期时不会被使用。

\textbf{边界实验。} 先用全奇数、全偶数、最大值变奇数停止、数值范围做反例检查。实验时覆盖全奇数、全偶数、最大值变奇数停止、数值范围，记录堆顶是否过期、比较器是否让正确候选在堆顶、K 或窗口大小为边界值时是否稳定。

\textbf{复杂度变化。} 暴力每轮扫描或排序动态候选，会把线性扫描或排序反复发生；堆把取极值、插入和删除边界控制在对数级。

\textbf{迁移追问。} 如果题目改成“这是堆维护动态极值和状态空间规范化的结合”，先判断原状态是否仍然覆盖新答案集合，再决定是微调边界、改 value 语义，还是换算法。

\noindent\textbf{LeetCode 1851 包含每个查询的最小区间。}

\textbf{题目说明。} LeetCode 1851 包含每个查询的最小区间 的题面入口要先还原成具体任务：每个查询点需要找覆盖它的最短区间，区间和查询都可离线排序。

\textbf{样例入口。} 拿 intervals=[[1, 4], [2, 4], [3, 6]]、queries=[2, 3, 5] 手算，查询 2 时可选 [1, 4] 和 [2, 4]，较短的是 [2, 4]。

\textbf{暴力基线。} 慢版本每轮扫描全部候选来找最大或最小，再更新候选集合。它的答案选择过程清楚，但动态集合每变化一次就重新找极值，代价太高。放到本题就是：先按“每个查询点需要找覆盖它的最短区间，区间和查询都可离线排序”完整试慢版本；样例里已经能看到“规律是离线按查询升序加入左端不大于查询的区间，堆顶清理右端已过期区间”，所以优化前必须先能说清慢版本枚举了哪些候选。

\textbf{暴力过程。} 拿 intervals=[[1, 4], [2, 4], [3, 6]]、queries=[2, 3, 5] 手算，查询 2 时可选 [1, 4] 和 [2, 4]，较短的是 [2, 4]。

\textbf{重复工作。} 题面模型：每个查询点需要找覆盖它的最短区间，区间和查询都可离线排序。暴力版的慢点要落到本题：慢版本会围绕“每个查询点需要找覆盖它的最短区间，区间和查询都可离线排序”反复扫描动态候选集找极值。样例规律是“规律是离线按查询升序加入左端不大于查询的区间，堆顶清理右端已过期区间”，说明每轮真正需要的是堆顶边界，而不是把候选全集重新排序。后面的优化只消掉这类重复工作，不能改变答案集合。

\textbf{优化条件。} 本题的关键观察是：按查询值升序扫描，把左端不大于查询的区间入堆，堆顶过期区间弹出。这句话必须能翻译成可维护状态；如果题目条件改变后观察不再成立，就不能继续套原来的写法。

\textbf{相邻状态。} 规律是离线按查询升序加入左端不大于查询的区间，堆顶清理右端已过期区间。把这个特例继续拆开看：堆顶必须正好代表下一步最需要处理的候选；若堆里可能有过期对象，读取堆顶前要先清理。堆只解决动态极值，不自动证明被弹出或被丢弃的元素无用。

\textbf{状态复用。} 把每轮全量扫描改成堆顶查询；插入、弹出和过期清理共同维护动态候选集。本题落点是：规律是离线按查询升序加入左端不大于查询的区间，堆顶清理右端已过期区间。手算时只改被新输入影响的那一小部分，而不是回头重算完整候选。

\textbf{指针和字段。} 状态字段要能说清：堆元素字段、堆顶比较规则、候选是否过期、答案读取时机；如果需要懒删除，还要保存删除计数。手算时按这个协议跟踪：拿 intervals=[[1, 4], [2, 4], [3, 6]]、queries=[2, 3, 5] 手算，查询 2 时可选 [1, 4] 和 [2, 4]，较短的是 [2, 4]。然后按协议复查：手算时记录堆顶、候选集合和是否有过期元素。每次使用堆顶前都要确认它仍然属于当前问题。

\textbf{代码落点。} 比较器方向先用小样例验证；每次使用堆顶前清理过期项；堆中保存下标时要能回查原数据。关键代码行必须能逐句翻译成“旧状态怎样变成新状态”，否则就只是模板。

\textbf{正确性。} 证明从候选覆盖开始：所有可能成为下一步答案的对象都在堆中，堆顶过期时不会被使用。

\textbf{边界实验。} 先用没有覆盖区间、区间长度相同、查询重复、右端过期做反例检查。实验时覆盖没有覆盖区间、区间长度相同、查询重复、右端过期，记录堆顶是否过期、比较器是否让正确候选在堆顶、K 或窗口大小为边界值时是否稳定。

\textbf{复杂度变化。} 暴力每轮扫描或排序动态候选，会把线性扫描或排序反复发生；堆把取极值、插入和删除边界控制在对数级。

\textbf{迁移追问。} 如果题目改成“这是离线扫描线、堆和区间覆盖的组合题”，先判断原状态是否仍然覆盖新答案集合，再决定是微调边界、改 value 语义，还是换算法。

\subsection{实验复盘}

追问通常会问为什么不完整排序、堆顶是否可能过期、比较器方向是否写反。请准备说明堆里元素的业务含义，而不是只说 priority\_queue。

复盘本节时先从 LeetCode 23 合并 K 个升序链表、LeetCode 215 数组中的第 K 个最大元素、LeetCode 295 数据流的中位数 开始：每道题都先手写一个能覆盖全部候选的暴力版，再指出优化结构具体保存了哪一段历史信息，最后用相邻案例比较为什么不能互换解法。

贪心证明深度题卡 按题型拆成章内大节，但每个大节都先从 LeetCode 案例切入，再进入分流、案例矩阵、案例推导和实验复盘。这样读者看到的是从具体题推到规律，而不是先读抽象方法再猜它能用在哪。

\topic{贪心与交换证明}
这一节先看 LeetCode 45 跳跃游戏 II、LeetCode 55 跳跃游戏、LeetCode 122 买卖股票 II 这几道 LeetCode 题，再整理同一题型的分流和推导。阅读顺序不是先背原理，而是先看案例的暴力瓶颈，再把重复出现的观察抽象成方法。

\subsection{原理和适用条件}

以 LeetCode 45 跳跃游戏 II、LeetCode 55 跳跃游戏、LeetCode 122 买卖股票 II 为主线：LeetCode 45 跳跃游戏 II：模型是“每一次跳跃可以看成 BFS 的一层，当前层内所有位置共享同一次跳数。”，优化抓手是“扫描当前层时维护下一层最远边界，到达当前层末尾才增加跳数。”；LeetCode 55 跳跃游戏：模型是“扫描过程中只关心当前能到达的最远位置。”，优化抓手是“如果下标超过最远可达，任何之前路径都不能到达它；否则用它扩展边界。”；LeetCode 122 买卖股票 II：模型是“允许多次交易且不能同时持有，所有正收益差都可以收集。”，优化抓手是“把长上涨段拆成每天买卖收益相同，因此累加正差值最优。”。这些案例共同推出的规律是：先把选择空间完整枚举，再寻找可以交换到最优解前面的局部选择。贪心必须能证明当前选择不会让后续资源变差，而不是只看排序后代码短。 方法名只是复盘后的标签，真正要掌握的是从题面约束、暴力失败、状态设计到边界反例的推导链。

\begin{principlebox}{图示：从 LeetCode 案例到 贪心与交换证明推导链}
\begin{tabular}{@{}p{0.18\linewidth}p{0.74\linewidth}@{}}
暴力基线 & 枚举候选、完整模拟或逐个检查，先保证覆盖全部可能答案。\\
瓶颈定位 & 慢点是选择空间爆炸。若每一步都有很多候选，而某个排序规则或边界规则能安全丢掉大部分候选，就可能存在贪心。\\
结构选择 & 贪心需要找到可交换的局部最优：最早结束、最小右端、当前最远覆盖、当前最大收益或最小代价。排序只是手段，不是证明。\\
不变量 & 贪心不变量通常是当前方案在相同选择次数下不落后于任何其他方案，或者当前保留的资源最多、结束最早、右边界最小。\\
代码落地 & 实现时先把比较器写成明确的业务顺序。区间题注意起点相同时终点升序还是降序；覆盖题注意闭区间和半开区间的等号。\\
\end{tabular}
\end{principlebox}

\subsection{LeetCode 案例分流}

\begin{principlebox}{从 LeetCode 案例分流：贪心与交换证明}
\textbf{最早结束和区间选择。}
代表案例：LeetCode 435 无重叠区间、LeetCode 452 引爆气球、LeetCode 646 最长数对链。先看这些题的题面条件：选择一个对象会占用时间或空间区间。由此推出的处理方式是：按结束位置排序，优先保留结束最早的候选。风险点：必须用交换论证，不是所有按端点排序都正确。

\textbf{最远覆盖。}
代表案例：LeetCode 55 跳跃游戏、LeetCode 45 跳跃游戏 II、LeetCode 763 划分字母区间。先看这些题的题面条件：当前选择决定下一段可覆盖范围。由此推出的处理方式是：扫描时维护当前可达最远边界或当前片段右端。风险点：当前层边界和下一层最远边界不能混。

\textbf{局部约束两遍扫描。}
代表案例：LeetCode 135 分发糖果、LeetCode 42 接雨水前后缀思路。先看这些题的题面条件：相邻约束可从左右两个方向分别满足。由此推出的处理方式是：左到右满足一侧约束，右到左满足另一侧，最后合并。风险点：只扫一遍往往不能同时满足两侧关系。

\textbf{资源反悔。}
代表案例：LeetCode 630 课程表 III、LeetCode 871 最低加油次数。先看这些题的题面条件：当前选了某些对象后，未来遇到更优对象可能需要替换。由此推出的处理方式是：先接受候选，用堆或排序在超约束时丢掉最差选择。风险点：反悔贪心要证明丢掉的候选不会让结果变差。

\end{principlebox}

\subsection{案例矩阵}

本节案例覆盖 LeetCode 45 跳跃游戏 II、LeetCode 55 跳跃游戏、LeetCode 122 买卖股票 II、LeetCode 134 加油站、LeetCode 135 分发糖果、LeetCode 169 多数元素、LeetCode 316 去除重复字母、LeetCode 406 根据身高重建队列、LeetCode 435 无重叠区间、LeetCode 452 用最少数量的箭引爆气球、LeetCode 630 课程表 III、LeetCode 646 最长数对链、LeetCode 678 有效的括号字符串、LeetCode 763 划分字母区间、LeetCode 860 柠檬水找零。阅读时不要按题号背答案，而要先判断题面落在哪个分流场景：查询目标是什么，输入约束给了什么单调性或等价关系，暴力慢在哪里，优化结构保存了什么状态。

\begin{principlebox}{案例矩阵}
\textbf{LeetCode 45 跳跃游戏 II。}
每一次跳跃可以看成 BFS 的一层，当前层内所有位置共享同一次跳数。单点风险：数组长度一、刚好到达、存在多个可选跳点、不要在每个位置都加跳数。

\textbf{LeetCode 55 跳跃游戏。}
扫描过程中只关心当前能到达的最远位置。单点风险：第一个元素为零、数组长度一、恰好到达最后、远超最后。

\textbf{LeetCode 122 买卖股票 II。}
允许多次交易且不能同时持有，所有正收益差都可以收集。单点风险：下降段、平台价格、只有一天、手续费不存在。

\textbf{LeetCode 134 加油站。}
绕行一圈是否可行取决于总油量差和局部剩余油量。单点风险：总和刚好为零、多个可行起点、某站油量为零、消耗大于补给。

\textbf{LeetCode 135 分发糖果。}
相邻评分约束需要同时满足左邻和右邻两个方向。单点风险：评分相等、严格递增、严格递减、山峰和谷底。

\textbf{LeetCode 169 多数元素。}
超过一半的元素可以和其他元素两两抵消后仍然剩下。单点风险：题目是否保证存在多数、计数归零、全相同。

\textbf{LeetCode 316 去除重复字母。}
每个字符必须保留一次，同时结果字典序尽量小。单点风险：字符只出现一次、重复字符、弹出后 visited 恢复、字典序比较。

\textbf{LeetCode 406 根据身高重建队列。}
高个子的位置先确定后，不会被矮个子影响其前方高个数量。单点风险：相同身高、k 为零、插入到末尾、结果稳定性。

\textbf{LeetCode 435 无重叠区间。}
保留最多不重叠区间等价于删除最少区间。单点风险：端点相接是否重叠、完全覆盖、相同终点、空列表。

\textbf{LeetCode 452 用最少数量的箭引爆气球。}
一支箭的位置可以覆盖所有包含该位置的区间。单点风险：端点相同、负坐标、区间包含、闭区间边界。

\textbf{LeetCode 630 课程表 III。}
每门课有持续时间和截止日，已选课程总时长不能超过当前截止日。单点风险：课程无法单独完成、截止日相同、替换已选课程、空列表。

\textbf{LeetCode 646 最长数对链。}
每个数对像一个区间，选择链要求前一个右端小于后一个左端。单点风险：端点相等不允许连接、完全覆盖、负数端点、空数组。

\textbf{LeetCode 678 有效的括号字符串。}
星号可以作为左括号、右括号或空，使未匹配左括号数量变成一个可能区间。单点风险：上界变负、最终下界不为零、连续星号、前缀右括号过多。

\textbf{LeetCode 763 划分字母区间。}
每个片段必须覆盖其中所有字符的最后出现位置。单点风险：单字符、所有字符相同、字符多次跨段。

\textbf{LeetCode 860 柠檬水找零。}
找零资源只有 5 元和 10 元，5 元比 10 元更灵活。单点风险：第一张不是 5、连续 10、连续 20、只有 5、5 元数量不足。

\end{principlebox}

\subsection{共性落点}

\begin{principlebox}{本组共性落点}
\textbf{慢版本。} 暴力做法枚举所有选择顺序或所有子集，找出最优方案。贪心题的难点不在写循环，而在证明局部选择不会破坏全局最优。
\textbf{瓶颈。} 慢点是选择空间爆炸。若每一步都有很多候选，而某个排序规则或边界规则能安全丢掉大部分候选，就可能存在贪心。
\textbf{状态字段。} 处理顺序、当前已选方案、剩余资源或最远边界、答案计数；若使用排序，要说明排序键；每次局部选择后要能说明当前状态不落后。
\textbf{不变量。} 贪心不变量通常是当前方案在相同选择次数下不落后于任何其他方案，或者当前保留的资源最多、结束最早、右边界最小。
\textbf{实现检查。} 把局部规则写成业务含义；若需要排序，就处理相同关键字；循环中只维护证明需要的状态，不把局部选择和答案更新混在一起。
\textbf{C++ 落地。} 实现时先把比较器写成明确的业务顺序。区间题注意起点相同时终点升序还是降序；覆盖题注意闭区间和半开区间的等号。
\textbf{证明抓手。} 证明从交换或领先性开始：任意最优解都能替换成当前贪心选择而不变差，或贪心状态始终不落后。
\end{principlebox}

\subsection{案例推导}

\noindent\textbf{LeetCode 45 跳跃游戏 II。}

\textbf{题目说明。} LeetCode 45 跳跃游戏 II 的题面入口要先还原成具体任务：每一次跳跃可以看成 BFS 的一层，当前层内所有位置共享同一次跳数。

\textbf{样例入口。} 拿 [2, 3, 1, 1, 4] 手算，起点这一层能到下标 1 和 2；扫描这一层时，下一层最远可到 4，走完当前层末尾才把步数加一。

\textbf{暴力基线。} 暴力 BFS 可以从每个当前下标扩展所有可跳位置，按层找第一次到终点的步数。

\textbf{暴力过程。} 以 [2, 3, 1, 1, 4] 为例，第一层是从 0 能到的 1、2；扫描这一层时，下一层最远可到 4，所以答案是 2。

\textbf{重复工作。} 题面模型：每一次跳跃可以看成 BFS 的一层，当前层内所有位置共享同一次跳数。暴力版的慢点要落到本题：浪费在显式保存一层内所有节点。数组下标的可达层是连续区间，只需要当前层右边界和下一层最远边界。后面的优化只消掉这类重复工作，不能改变答案集合。

\textbf{优化条件。} 本题的关键观察是：扫描当前层时维护下一层最远边界，到达当前层末尾才增加跳数。这句话必须能翻译成可维护状态；如果题目条件改变后观察不再成立，就不能继续套原来的写法。

\textbf{相邻状态。} 规律是把可达区间看成 BFS 层，当前层结束时切到下一层；不能在每个位置都加跳数。把这个特例继续拆开看：先构造一个非贪心选择，再比较两者留下的资源、右边界或剩余时间。只有能把非贪心第一步交换成贪心第一步且不变差，局部选择才是规律。

\textbf{状态复用。} current\_end 表示当前跳数能覆盖的最右位置，farthest 表示下一跳能覆盖的最远位置。扫描到 current\_end 时，步数加一并切换到下一层。手算时只改被新输入影响的那一小部分，而不是回头重算完整候选。

\textbf{指针和字段。} 状态字段要能说清：i 是扫描位置，jumps 是已完成层数，current\_end 是本层末尾，farthest 是下一层末尾。手算时按这个协议跟踪：i=0 时 farthest=2，走到 current\_end=0 后 jumps=1、current\_end=2；扫描 i=1 时 farthest 更新到 4，走到 i=2 后 jumps=2。

\textbf{代码落点。} 关键是循环通常只扫到 n-2，因为到达最后一个位置不需要再跳。不能在每个位置都 jumps++，只能在层结束时加。关键代码行必须能逐句翻译成“旧状态怎样变成新状态”，否则就只是模板。

\textbf{正确性。} 所有当前跳数可达的位置构成一个连续区间；在这个区间内选择能让下一层覆盖最远的位置，不会比保存全部 BFS 节点差。

\textbf{边界实验。} 先用数组长度一、刚好到达、存在多个可选跳点、不要在每个位置都加跳数做反例检查。实验时除了数组长度一、刚好到达、存在多个可选跳点、不要在每个位置都加跳数，还要故意替换成一个看似合理但错误的局部规则，构造反例。能解释错误贪心为何失败，才算理解正确贪心。

\textbf{复杂度变化。} 显式 BFS 分支多；区间层扫描 O(n) 时间、O(1) 空间。

\textbf{迁移追问。} 如果题目改成“若问能否到达，只需维护全局最远可达位置”，先判断原状态是否仍然覆盖新答案集合，再决定是微调边界、改 value 语义，还是换算法。

\noindent\textbf{LeetCode 55 跳跃游戏。}

\textbf{题目说明。} LeetCode 55 跳跃游戏 的题面入口要先还原成具体任务：扫描过程中只关心当前能到达的最远位置。

\textbf{样例入口。} 拿 [2, 3, 1, 1, 4] 手算，下标 0 能覆盖到 2；走到下标 1 时最远可达更新到 4，已经覆盖终点。再拿 [0, 1]，扫描到下标 1 时已经超过最远可达，失败。

\textbf{暴力基线。} 暴力把每个下标能跳到的位置都当成分支搜索，尝试是否存在一条路径能到终点。

\textbf{暴力过程。} 以 [2, 3, 1, 1, 4] 为例，从 0 可以跳到 1 或 2；若到 1，最远能继续覆盖到 4。以 [0, 1] 为例，从 0 一步也走不出去。

\textbf{重复工作。} 题面模型：扫描过程中只关心当前能到达的最远位置。暴力版的慢点要落到本题：浪费在枚举具体路径。判断能否到达终点时，不关心走的是哪条路径，只关心目前所有可达位置能覆盖到的最远下标。后面的优化只消掉这类重复工作，不能改变答案集合。

\textbf{优化条件。} 本题的关键观察是：如果下标超过最远可达，任何之前路径都不能到达它；否则用它扩展边界。这句话必须能翻译成可维护状态；如果题目条件改变后观察不再成立，就不能继续套原来的写法。

\textbf{相邻状态。} 规律是只维护当前能到达的最远边界，任何超过边界的位置都不可能由前面路径到达。把这个特例继续拆开看：先构造一个非贪心选择，再比较两者留下的资源、右边界或剩余时间。只有能把非贪心第一步交换成贪心第一步且不变差，局部选择才是规律。

\textbf{状态复用。} 扫描下标 i，维护 farthest。若 i>farthest，说明当前下标不可达，直接失败；否则用 i+nums[i] 更新 farthest。手算时只改被新输入影响的那一小部分，而不是回头重算完整候选。

\textbf{指针和字段。} 状态字段要能说清：i 是当前检查位置，farthest 是所有已可达位置能跳到的最远边界，n-1 是目标边界。手算时按这个协议跟踪：样例中 i=0 后 farthest=2；i=1 可达并把 farthest 更新为 4，已经覆盖终点。

\textbf{代码落点。} 关键顺序是先判断当前位置是否超过 farthest；若没超过，再用 i+nums[i] 扩展最远可达位置。数组长度为 1 时起点就是终点。关键代码行必须能逐句翻译成“旧状态怎样变成新状态”，否则就只是模板。

\textbf{正确性。} 在扫描到 i 前，farthest 汇总了所有可达下标的最远覆盖；若 i 超过它，任何旧路径都到不了 i，更不可能通过 i 之后的位置到终点。

\textbf{边界实验。} 先用第一个元素为零、数组长度一、恰好到达最后、远超最后做反例检查。实验时除了第一个元素为零、数组长度一、恰好到达最后、远超最后，还要故意替换成一个看似合理但错误的局部规则，构造反例。能解释错误贪心为何失败，才算理解正确贪心。

\textbf{复杂度变化。} 路径搜索最坏指数级；贪心边界扫描 O(n) 时间、O(1) 空间。

\textbf{迁移追问。} 如果题目改成“求最少跳数时把可达区间看成 BFS 层”，先判断原状态是否仍然覆盖新答案集合，再决定是微调边界、改 value 语义，还是换算法。

\noindent\textbf{LeetCode 122 买卖股票 II。}

\textbf{题目说明。} LeetCode 122 买卖股票 II 的题面入口要先还原成具体任务：允许多次交易且不能同时持有，所有正收益差都可以收集。

\textbf{样例入口。} 拿 [7, 1, 5, 3, 6, 4] 手算，1 到 5 的利润 4，加上 3 到 6 的利润 3，总利润 7。把 1 到 6 的上涨段拆成每天正差值，收益相同。

\textbf{暴力基线。} 慢版本枚举选择顺序或用 DP 保留多个可能方案，先确认最优值存在。随后才检查局部选择是否能通过交换或领先性固定下来。放到本题就是：先按“允许多次交易且不能同时持有，所有正收益差都可以收集”完整试慢版本；样例里已经能看到“规律是无限次交易且无手续费时，所有正差值都可收集；下降段和平台不贡献收益”，所以优化前必须先能说清慢版本枚举了哪些候选。

\textbf{暴力过程。} 拿 [7, 1, 5, 3, 6, 4] 手算，1 到 5 的利润 4，加上 3 到 6 的利润 3，总利润 7。把 1 到 6 的上涨段拆成每天正差值，收益相同。

\textbf{重复工作。} 题面模型：允许多次交易且不能同时持有，所有正收益差都可以收集。暴力版的慢点要落到本题：慢版本会围绕“允许多次交易且不能同时持有，所有正收益差都可以收集”枚举选择顺序。样例规律是“规律是无限次交易且无手续费时，所有正差值都可收集；下降段和平台不贡献收益”，说明存在一个局部选择或边界状态可以安全保留；不能证明这个选择不变差时，就不能把它叫贪心。后面的优化只消掉这类重复工作，不能改变答案集合。

\textbf{优化条件。} 本题的关键观察是：把长上涨段拆成每天买卖收益相同，因此累加正差值最优。这句话必须能翻译成可维护状态；如果题目条件改变后观察不再成立，就不能继续套原来的写法。

\textbf{相邻状态。} 规律是无限次交易且无手续费时，所有正差值都可收集；下降段和平台不贡献收益。把这个特例继续拆开看：先构造一个非贪心选择，再比较两者留下的资源、右边界或剩余时间。只有能把非贪心第一步交换成贪心第一步且不变差，局部选择才是规律。

\textbf{状态复用。} 把指数级选择压成一个可证明安全的局部选择；状态通常是当前最远边界、最早结束时间、最小代价或剩余资源。本题落点是：规律是无限次交易且无手续费时，所有正差值都可收集；下降段和平台不贡献收益。手算时只改被新输入影响的那一小部分，而不是回头重算完整候选。

\textbf{指针和字段。} 状态字段要能说清：处理顺序、当前已选方案、剩余资源或最远边界、答案计数；若使用排序，要说明排序键；每次局部选择后要能说明当前状态不落后。手算时按这个协议跟踪：拿 [7, 1, 5, 3, 6, 4] 手算，1 到 5 的利润 4，加上 3 到 6 的利润 3，总利润 7。把 1 到 6 的上涨段拆成每天正差值，收益相同。然后按协议复查：手算时同时写一个非贪心选择，与贪心选择比较剩余资源。若无法说明贪心后剩余资源不差，就不能直接下结论。

\textbf{代码落点。} 把局部规则写成业务含义；若需要排序，就处理相同关键字；循环中只维护证明需要的状态，不把局部选择和答案更新混在一起。关键代码行必须能逐句翻译成“旧状态怎样变成新状态”，否则就只是模板。

\textbf{正确性。} 证明从交换或领先性开始：任意最优解都能替换成当前贪心选择而不变差，或贪心状态始终不落后。

\textbf{边界实验。} 先用下降段、平台价格、只有一天、手续费不存在做反例检查。实验时除了下降段、平台价格、只有一天、手续费不存在，还要故意替换成一个看似合理但错误的局部规则，构造反例。能解释错误贪心为何失败，才算理解正确贪心。

\textbf{复杂度变化。} 暴力枚举选择顺序可能指数级；贪心通常先排序，再线性扫描或用堆维护少量候选。

\textbf{迁移追问。} 如果题目改成“有手续费时不能简单累加正差，需要状态机或调整贪心阈值”，先判断原状态是否仍然覆盖新答案集合，再决定是微调边界、改 value 语义，还是换算法。

\noindent\textbf{LeetCode 134 加油站。}

\textbf{题目说明。} LeetCode 134 加油站 的题面入口要先还原成具体任务：绕行一圈是否可行取决于总油量差和局部剩余油量。

\textbf{样例入口。} 拿 gas=[1, 2, 3, 4, 5]、cost=[3, 4, 5, 1, 2] 手算，从 0 出发中途油量为负，0 到失败点之间任何站都不能作为起点；从 3 出发可以绕完。

\textbf{暴力基线。} 暴力把每个加油站都当起点，模拟绕一圈，若油量中途为负就换下一个起点。

\textbf{暴力过程。} gas=[1, 2, 3, 4, 5]、cost=[3, 4, 5, 1, 2] 中，从 0 出发会在中途油量为负；从 3 出发可以完成一圈。

\textbf{重复工作。} 题面模型：绕行一圈是否可行取决于总油量差和局部剩余油量。暴力版的慢点要落到本题：浪费在失败起点之后仍逐个重试中间站。若从 start 到 i 的累计油量已经为负，那么 start..i 中任何站都不可能作为可行起点。后面的优化只消掉这类重复工作，不能改变答案集合。

\textbf{优化条件。} 本题的关键观察是：总和为负必不可行；扫描时若当前剩余为负，起点只能放到下一站。这句话必须能翻译成可维护状态；如果题目条件改变后观察不再成立，就不能继续套原来的写法。

\textbf{相邻状态。} 规律是总油量差为负则无解，局部剩余为负时起点跳到下一站。把这个特例继续拆开看：先构造一个非贪心选择，再比较两者留下的资源、右边界或剩余时间。只有能把非贪心第一步交换成贪心第一步且不变差，局部选择才是规律。

\textbf{状态复用。} 扫描一次，total 记录总油量差，tank 记录当前候选起点到当前位置的剩余油量，start 是当前候选起点。手算时只改被新输入影响的那一小部分，而不是回头重算完整候选。

\textbf{指针和字段。} 状态字段要能说清：diff=gas[i]-cost[i]，tank 是局部剩余，total 是全局可行性，start 是下一次尝试起点。手算时按这个协议跟踪：样例前几站累计 tank 变负后，把 start 移到下一站并清零 tank；最后 total>=0，返回最终 start=3。

\textbf{代码落点。} 关键是 total 全程累加，tank<0 时 start=i+1、tank=0。循环结束后 total<0 返回 -1，否则返回 start。关键代码行必须能逐句翻译成“旧状态怎样变成新状态”，否则就只是模板。

\textbf{正确性。} tank<0 说明从当前 start 到 i 的总补给不足；其中任意中间站作为起点，少了前面一段补给，只会更差，故可整体跳过。

\textbf{边界实验。} 先用总和刚好为零、多个可行起点、某站油量为零、消耗大于补给做反例检查。实验时除了总和刚好为零、多个可行起点、某站油量为零、消耗大于补给，还要故意替换成一个看似合理但错误的局部规则，构造反例。能解释错误贪心为何失败，才算理解正确贪心。

\textbf{复杂度变化。} 暴力 O(\ensuremath{n^{2}})，一次扫描 O(n) 时间、O(1) 空间。

\textbf{迁移追问。} 如果题目改成“它的领先性证明是前一段任意站点都不能作为起点”，先判断原状态是否仍然覆盖新答案集合，再决定是微调边界、改 value 语义，还是换算法。

\noindent\textbf{LeetCode 135 分发糖果。}

\textbf{题目说明。} LeetCode 135 分发糖果 的题面入口要先还原成具体任务：相邻评分约束需要同时满足左邻和右邻两个方向。

\textbf{样例入口。} 拿 ratings=[1, 0, 2] 手算，中间孩子评分最低，左右都要比它多糖，结果 [2, 1, 2]。只从左到右会漏掉右侧约束。

\textbf{暴力基线。} 慢版本枚举选择顺序或用 DP 保留多个可能方案，先确认最优值存在。随后才检查局部选择是否能通过交换或领先性固定下来。放到本题就是：先按“相邻评分约束需要同时满足左邻和右邻两个方向”完整试慢版本；样例里已经能看到“规律是左到右满足左邻递增，右到左满足右邻递增，最终取两侧要求最大值”，所以优化前必须先能说清慢版本枚举了哪些候选。

\textbf{暴力过程。} 拿 ratings=[1, 0, 2] 手算，中间孩子评分最低，左右都要比它多糖，结果 [2, 1, 2]。只从左到右会漏掉右侧约束。

\textbf{重复工作。} 题面模型：相邻评分约束需要同时满足左邻和右邻两个方向。暴力版的慢点要落到本题：慢版本会围绕“相邻评分约束需要同时满足左邻和右邻两个方向”枚举选择顺序。样例规律是“规律是左到右满足左邻递增，右到左满足右邻递增，最终取两侧要求最大值”，说明存在一个局部选择或边界状态可以安全保留；不能证明这个选择不变差时，就不能把它叫贪心。后面的优化只消掉这类重复工作，不能改变答案集合。

\textbf{优化条件。} 本题的关键观察是：左到右满足递增关系，右到左满足反向递增关系，最后取两侧要求最大值。这句话必须能翻译成可维护状态；如果题目条件改变后观察不再成立，就不能继续套原来的写法。

\textbf{相邻状态。} 规律是左到右满足左邻递增，右到左满足右邻递增，最终取两侧要求最大值。把这个特例继续拆开看：先构造一个非贪心选择，再比较两者留下的资源、右边界或剩余时间。只有能把非贪心第一步交换成贪心第一步且不变差，局部选择才是规律。

\textbf{状态复用。} 把指数级选择压成一个可证明安全的局部选择；状态通常是当前最远边界、最早结束时间、最小代价或剩余资源。本题落点是：规律是左到右满足左邻递增，右到左满足右邻递增，最终取两侧要求最大值。手算时只改被新输入影响的那一小部分，而不是回头重算完整候选。

\textbf{指针和字段。} 状态字段要能说清：处理顺序、当前已选方案、剩余资源或最远边界、答案计数；若使用排序，要说明排序键；每次局部选择后要能说明当前状态不落后。手算时按这个协议跟踪：拿 ratings=[1, 0, 2] 手算，中间孩子评分最低，左右都要比它多糖，结果 [2, 1, 2]。只从左到右会漏掉右侧约束。然后按协议复查：手算时同时写一个非贪心选择，与贪心选择比较剩余资源。若无法说明贪心后剩余资源不差，就不能直接下结论。

\textbf{代码落点。} 把局部规则写成业务含义；若需要排序，就处理相同关键字；循环中只维护证明需要的状态，不把局部选择和答案更新混在一起。关键代码行必须能逐句翻译成“旧状态怎样变成新状态”，否则就只是模板。

\textbf{正确性。} 证明从交换或领先性开始：任意最优解都能替换成当前贪心选择而不变差，或贪心状态始终不落后。

\textbf{边界实验。} 先用评分相等、严格递增、严格递减、山峰和谷底做反例检查。实验时除了评分相等、严格递增、严格递减、山峰和谷底，还要故意替换成一个看似合理但错误的局部规则，构造反例。能解释错误贪心为何失败，才算理解正确贪心。

\textbf{复杂度变化。} 暴力枚举选择顺序可能指数级；贪心通常先排序，再线性扫描或用堆维护少量候选。

\textbf{迁移追问。} 如果题目改成“这是局部约束两遍扫描，不能只用一个方向贪心”，先判断原状态是否仍然覆盖新答案集合，再决定是微调边界、改 value 语义，还是换算法。

\noindent\textbf{LeetCode 169 多数元素。}

\textbf{题目说明。} LeetCode 169 多数元素 的题面入口要先还原成具体任务：超过一半的元素可以和其他元素两两抵消后仍然剩下。

\textbf{样例入口。} 拿 [2, 2, 1, 1, 1, 2, 2] 手算，候选和不同元素互相抵消，最后剩下的候选是 2。

\textbf{暴力基线。} 慢版本枚举选择顺序或用 DP 保留多个可能方案，先确认最优值存在。随后才检查局部选择是否能通过交换或领先性固定下来。放到本题就是：先按“超过一半的元素可以和其他元素两两抵消后仍然剩下”完整试慢版本；样例里已经能看到“规律是多数元素超过一半，和其他元素两两抵消后仍会剩下；题目不保证存在时需要二次计数验证”，所以优化前必须先能说清慢版本枚举了哪些候选。

\textbf{暴力过程。} 拿 [2, 2, 1, 1, 1, 2, 2] 手算，候选和不同元素互相抵消，最后剩下的候选是 2。

\textbf{重复工作。} 题面模型：超过一半的元素可以和其他元素两两抵消后仍然剩下。暴力版的慢点要落到本题：慢版本会围绕“超过一半的元素可以和其他元素两两抵消后仍然剩下”枚举选择顺序。样例规律是“规律是多数元素超过一半，和其他元素两两抵消后仍会剩下；题目不保证存在时需要二次计数验证”，说明存在一个局部选择或边界状态可以安全保留；不能证明这个选择不变差时，就不能把它叫贪心。后面的优化只消掉这类重复工作，不能改变答案集合。

\textbf{优化条件。} 本题的关键观察是：Boyer-Moore 投票维护候选和计数。这句话必须能翻译成可维护状态；如果题目条件改变后观察不再成立，就不能继续套原来的写法。

\textbf{相邻状态。} 规律是多数元素超过一半，和其他元素两两抵消后仍会剩下；题目不保证存在时需要二次计数验证。把这个特例继续拆开看：先构造一个非贪心选择，再比较两者留下的资源、右边界或剩余时间。只有能把非贪心第一步交换成贪心第一步且不变差，局部选择才是规律。

\textbf{状态复用。} 把指数级选择压成一个可证明安全的局部选择；状态通常是当前最远边界、最早结束时间、最小代价或剩余资源。本题落点是：规律是多数元素超过一半，和其他元素两两抵消后仍会剩下；题目不保证存在时需要二次计数验证。手算时只改被新输入影响的那一小部分，而不是回头重算完整候选。

\textbf{指针和字段。} 状态字段要能说清：处理顺序、当前已选方案、剩余资源或最远边界、答案计数；若使用排序，要说明排序键；每次局部选择后要能说明当前状态不落后。手算时按这个协议跟踪：拿 [2, 2, 1, 1, 1, 2, 2] 手算，候选和不同元素互相抵消，最后剩下的候选是 2。然后按协议复查：手算时同时写一个非贪心选择，与贪心选择比较剩余资源。若无法说明贪心后剩余资源不差，就不能直接下结论。

\textbf{代码落点。} 把局部规则写成业务含义；若需要排序，就处理相同关键字；循环中只维护证明需要的状态，不把局部选择和答案更新混在一起。关键代码行必须能逐句翻译成“旧状态怎样变成新状态”，否则就只是模板。

\textbf{正确性。} 证明从交换或领先性开始：任意最优解都能替换成当前贪心选择而不变差，或贪心状态始终不落后。

\textbf{边界实验。} 先用题目是否保证存在多数、计数归零、全相同做反例检查。实验时除了题目是否保证存在多数、计数归零、全相同，还要故意替换成一个看似合理但错误的局部规则，构造反例。能解释错误贪心为何失败，才算理解正确贪心。

\textbf{复杂度变化。} 暴力枚举选择顺序可能指数级；贪心通常先排序，再线性扫描或用堆维护少量候选。

\textbf{迁移追问。} 如果题目改成“若找超过三分之一的元素，需要维护两个候选”，先判断原状态是否仍然覆盖新答案集合，再决定是微调边界、改 value 语义，还是换算法。

\noindent\textbf{LeetCode 316 去除重复字母。}

\textbf{题目说明。} LeetCode 316 去除重复字母 的题面入口要先还原成具体任务：每个字符必须保留一次，同时结果字典序尽量小。

\textbf{样例入口。} 拿 cbacdcbc 手算，遇到 b 时，前面的 c 后面还会出现且 c>b，所以可以弹出 c；但遇到只剩最后一次的字符就不能弹。

\textbf{暴力基线。} 慢版本枚举选择顺序或用 DP 保留多个可能方案，先确认最优值存在。随后才检查局部选择是否能通过交换或领先性固定下来。放到本题就是：先按“每个字符必须保留一次，同时结果字典序尽量小”完整试慢版本；样例里已经能看到“规律是单调栈维护结果字典序，弹出条件必须同时满足栈顶更大且后面还会出现”，所以优化前必须先能说清慢版本枚举了哪些候选。

\textbf{暴力过程。} 拿 cbacdcbc 手算，遇到 b 时，前面的 c 后面还会出现且 c>b，所以可以弹出 c；但遇到只剩最后一次的字符就不能弹。

\textbf{重复工作。} 题面模型：每个字符必须保留一次，同时结果字典序尽量小。暴力版的慢点要落到本题：慢版本会围绕“每个字符必须保留一次，同时结果字典序尽量小”枚举选择顺序。样例规律是“规律是单调栈维护结果字典序，弹出条件必须同时满足栈顶更大且后面还会出现”，说明存在一个局部选择或边界状态可以安全保留；不能证明这个选择不变差时，就不能把它叫贪心。后面的优化只消掉这类重复工作，不能改变答案集合。

\textbf{优化条件。} 本题的关键观察是：单调栈维护当前结果，若栈顶更大且后面还会出现，就可以弹出再等后面放回。这句话必须能翻译成可维护状态；如果题目条件改变后观察不再成立，就不能继续套原来的写法。

\textbf{相邻状态。} 规律是单调栈维护结果字典序，弹出条件必须同时满足栈顶更大且后面还会出现。把这个特例继续拆开看：先构造一个非贪心选择，再比较两者留下的资源、右边界或剩余时间。只有能把非贪心第一步交换成贪心第一步且不变差，局部选择才是规律。

\textbf{状态复用。} 把指数级选择压成一个可证明安全的局部选择；状态通常是当前最远边界、最早结束时间、最小代价或剩余资源。本题落点是：规律是单调栈维护结果字典序，弹出条件必须同时满足栈顶更大且后面还会出现。手算时只改被新输入影响的那一小部分，而不是回头重算完整候选。

\textbf{指针和字段。} 状态字段要能说清：处理顺序、当前已选方案、剩余资源或最远边界、答案计数；若使用排序，要说明排序键；每次局部选择后要能说明当前状态不落后。手算时按这个协议跟踪：拿 cbacdcbc 手算，遇到 b 时，前面的 c 后面还会出现且 c>b，所以可以弹出 c；但遇到只剩最后一次的字符就不能弹。然后按协议复查：手算时同时写一个非贪心选择，与贪心选择比较剩余资源。若无法说明贪心后剩余资源不差，就不能直接下结论。

\textbf{代码落点。} 把局部规则写成业务含义；若需要排序，就处理相同关键字；循环中只维护证明需要的状态，不把局部选择和答案更新混在一起。关键代码行必须能逐句翻译成“旧状态怎样变成新状态”，否则就只是模板。

\textbf{正确性。} 证明从交换或领先性开始：任意最优解都能替换成当前贪心选择而不变差，或贪心状态始终不落后。

\textbf{边界实验。} 先用字符只出现一次、重复字符、弹出后 visited 恢复、字典序比较做反例检查。实验时除了字符只出现一次、重复字符、弹出后 visited 恢复、字典序比较，还要故意替换成一个看似合理但错误的局部规则，构造反例。能解释错误贪心为何失败，才算理解正确贪心。

\textbf{复杂度变化。} 暴力枚举选择顺序可能指数级；贪心通常先排序，再线性扫描或用堆维护少量候选。

\textbf{迁移追问。} 如果题目改成“402 移掉 K 位数字也是单调栈删除，但删除预算和必须保留集合不同”，先判断原状态是否仍然覆盖新答案集合，再决定是微调边界、改 value 语义，还是换算法。

\noindent\textbf{LeetCode 406 根据身高重建队列。}

\textbf{题目说明。} LeetCode 406 根据身高重建队列 的题面入口要先还原成具体任务：高个子的位置先确定后，不会被矮个子影响其前方高个数量。

\textbf{样例入口。} 拿 [7, 0]、[7, 1]、[5, 0] 手算，先放高个子，矮个子插入不会改变高个子前面高个数量；[5, 0] 插到最前也不影响 [7, 0] 的统计。

\textbf{暴力基线。} 慢版本枚举选择顺序或用 DP 保留多个可能方案，先确认最优值存在。随后才检查局部选择是否能通过交换或领先性固定下来。放到本题就是：先按“高个子的位置先确定后，不会被矮个子影响其前方高个数量”完整试慢版本；样例里已经能看到“规律是按身高降序、k 升序排序，再按 k 插入”，所以优化前必须先能说清慢版本枚举了哪些候选。

\textbf{暴力过程。} 拿 [7, 0]、[7, 1]、[5, 0] 手算，先放高个子，矮个子插入不会改变高个子前面高个数量；[5, 0] 插到最前也不影响 [7, 0] 的统计。

\textbf{重复工作。} 题面模型：高个子的位置先确定后，不会被矮个子影响其前方高个数量。暴力版的慢点要落到本题：慢版本会围绕“高个子的位置先确定后，不会被矮个子影响其前方高个数量”枚举选择顺序。样例规律是“规律是按身高降序、k 升序排序，再按 k 插入”，说明存在一个局部选择或边界状态可以安全保留；不能证明这个选择不变差时，就不能把它叫贪心。后面的优化只消掉这类重复工作，不能改变答案集合。

\textbf{优化条件。} 本题的关键观察是：按身高降序、k 升序排序，再按 k 插入当前结果位置。这句话必须能翻译成可维护状态；如果题目条件改变后观察不再成立，就不能继续套原来的写法。

\textbf{相邻状态。} 规律是按身高降序、k 升序排序，再按 k 插入。把这个特例继续拆开看：先构造一个非贪心选择，再比较两者留下的资源、右边界或剩余时间。只有能把非贪心第一步交换成贪心第一步且不变差，局部选择才是规律。

\textbf{状态复用。} 把指数级选择压成一个可证明安全的局部选择；状态通常是当前最远边界、最早结束时间、最小代价或剩余资源。本题落点是：规律是按身高降序、k 升序排序，再按 k 插入。手算时只改被新输入影响的那一小部分，而不是回头重算完整候选。

\textbf{指针和字段。} 状态字段要能说清：处理顺序、当前已选方案、剩余资源或最远边界、答案计数；若使用排序，要说明排序键；每次局部选择后要能说明当前状态不落后。手算时按这个协议跟踪：拿 [7, 0]、[7, 1]、[5, 0] 手算，先放高个子，矮个子插入不会改变高个子前面高个数量；[5, 0] 插到最前也不影响 [7, 0] 的统计。然后按协议复查：手算时同时写一个非贪心选择，与贪心选择比较剩余资源。若无法说明贪心后剩余资源不差，就不能直接下结论。

\textbf{代码落点。} 把局部规则写成业务含义；若需要排序，就处理相同关键字；循环中只维护证明需要的状态，不把局部选择和答案更新混在一起。关键代码行必须能逐句翻译成“旧状态怎样变成新状态”，否则就只是模板。

\textbf{正确性。} 证明从交换或领先性开始：任意最优解都能替换成当前贪心选择而不变差，或贪心状态始终不落后。

\textbf{边界实验。} 先用相同身高、k 为零、插入到末尾、结果稳定性做反例检查。实验时除了相同身高、k 为零、插入到末尾、结果稳定性，还要故意替换成一个看似合理但错误的局部规则，构造反例。能解释错误贪心为何失败，才算理解正确贪心。

\textbf{复杂度变化。} 暴力枚举选择顺序可能指数级；贪心通常先排序，再线性扫描或用堆维护少量候选。

\textbf{迁移追问。} 如果题目改成“这是排序后插入的构造型贪心，证明依赖高个子先固定”，先判断原状态是否仍然覆盖新答案集合，再决定是微调边界、改 value 语义，还是换算法。

\noindent\textbf{LeetCode 435 无重叠区间。}

\textbf{题目说明。} LeetCode 435 无重叠区间 的题面入口要先还原成具体任务：保留最多不重叠区间等价于删除最少区间。

\textbf{样例入口。} 拿 [[1, 2], [2, 3], [3, 4], [1, 3]] 手算，选 [1, 2] 后还能接 [2, 3] 和 [3, 4]；若先选 [1, 3]，会挡住更多区间。

\textbf{暴力基线。} 暴力可以枚举保留哪些区间，再检查它们是否互不重叠，最后用总数减最大保留数量。

\textbf{暴力过程。} 以 [[1, 2], [2, 3], [3, 4], [1, 3]] 为例，选 [1, 2] 后还能接 [2, 3] 和 [3, 4]；若先选 [1, 3] 会挡住更多后续区间。

\textbf{重复工作。} 题面模型：保留最多不重叠区间等价于删除最少区间。暴力版的慢点要落到本题：浪费在枚举保留集合。区间选择里结束越早，给后面留下的空间越大，可以作为安全局部选择。后面的优化只消掉这类重复工作，不能改变答案集合。

\textbf{优化条件。} 本题的关键观察是：按终点升序选择，结束越早给后续留下空间越多。这句话必须能翻译成可维护状态；如果题目条件改变后观察不再成立，就不能继续套原来的写法。

\textbf{相邻状态。} 规律是按右端点升序选择，结束越早给未来留下的空间越多，删除数等于总数减保留数。把这个特例继续拆开看：先构造一个非贪心选择，再比较两者留下的资源、右边界或剩余时间。只有能把非贪心第一步交换成贪心第一步且不变差，局部选择才是规律。

\textbf{状态复用。} 按右端点升序排序。维护 last\_end；若当前区间 start>=last\_end，就保留它并更新 last\_end，否则删除它。手算时只改被新输入影响的那一小部分，而不是回头重算完整候选。

\textbf{指针和字段。} 状态字段要能说清：last\_end 是已保留区间的最后右端，kept 是保留数量，removed 是删除数量。手算时按这个协议跟踪：排序后先保留 [1, 2]；[1, 3] 与它重叠删除；[2, 3] 和 [3, 4] 都能接上。

\textbf{代码落点。} 关键是端点相接是否算重叠。本题中 [1, 2] 和 [2, 3] 不重叠，所以条件是 start>=last\_end。关键代码行必须能逐句翻译成“旧状态怎样变成新状态”，否则就只是模板。

\textbf{正确性。} 任意最优解的第一个区间都可替换为结束最早的区间，后续可用空间不变小；重复此过程得到贪心解。

\textbf{边界实验。} 先用端点相接是否重叠、完全覆盖、相同终点、空列表做反例检查。实验时除了端点相接是否重叠、完全覆盖、相同终点、空列表，还要故意替换成一个看似合理但错误的局部规则，构造反例。能解释错误贪心为何失败，才算理解正确贪心。

\textbf{复杂度变化。} 暴力指数级；排序 O(n log n)，扫描 O(n)，空间取决于排序。

\textbf{迁移追问。} 如果题目改成“用交换论证说明最早结束的选择不劣于其他第一选择”，先判断原状态是否仍然覆盖新答案集合，再决定是微调边界、改 value 语义，还是换算法。

\noindent\textbf{LeetCode 452 用最少数量的箭引爆气球。}

\textbf{题目说明。} LeetCode 452 用最少数量的箭引爆气球 的题面入口要先还原成具体任务：一支箭的位置可以覆盖所有包含该位置的区间。

\textbf{样例入口。} 拿 [[10, 16], [2, 8], [1, 6], [7, 12]] 手算，按右端排序后第一箭射在 6，可以覆盖 [1, 6] 和 [2, 8]；下一箭射在 12。

\textbf{暴力基线。} 暴力可以枚举箭的位置，或对每个气球单独射箭后再尝试合并覆盖关系。

\textbf{暴力过程。} 以 [[10, 16], [2, 8], [1, 6], [7, 12]] 为例，按右端排序后第一支箭射在 6，可同时覆盖 [1, 6] 和 [2, 8]。

\textbf{重复工作。} 题面模型：一支箭的位置可以覆盖所有包含该位置的区间。暴力版的慢点要落到本题：浪费在没有利用区间右端。第一支箭若要覆盖最早结束的气球，射在它右端能给后续重叠区间最大机会。后面的优化只消掉这类重复工作，不能改变答案集合。

\textbf{优化条件。} 本题的关键观察是：按右端点排序，当前箭放在最早结束气球的右端。这句话必须能翻译成可维护状态；如果题目条件改变后观察不再成立，就不能继续套原来的写法。

\textbf{相邻状态。} 规律是每次把箭放在最早结束气球的右端，能最大化后续空间。把这个特例继续拆开看：先构造一个非贪心选择，再比较两者留下的资源、右边界或剩余时间。只有能把非贪心第一步交换成贪心第一步且不变差，局部选择才是规律。

\textbf{状态复用。} 按 right 升序排序。arrow\_pos 初始为第一个区间右端；后续区间若 start<=arrow\_pos 就被当前箭覆盖，否则新增一支箭。手算时只改被新输入影响的那一小部分，而不是回头重算完整候选。

\textbf{指针和字段。} 状态字段要能说清：arrow\_pos 是当前箭的位置，arrows 是箭数，interval 是当前气球闭区间。手算时按这个协议跟踪：第一箭在 6，覆盖 [1, 6] 和 [2, 8]；遇到 [7, 12] 时 7>6，需要第二箭放到 12。

\textbf{代码落点。} 关键是闭区间边界用 start<=arrow\_pos 表示能覆盖。坐标可能为负，不影响排序和比较。关键代码行必须能逐句翻译成“旧状态怎样变成新状态”，否则就只是模板。

\textbf{正确性。} 最早结束区间必须被某支箭覆盖；把这支箭移动到它的右端不会减少已覆盖区间，还可能覆盖更多后续区间。

\textbf{边界实验。} 先用端点相同、负坐标、区间包含、闭区间边界做反例检查。实验时除了端点相同、负坐标、区间包含、闭区间边界，还要故意替换成一个看似合理但错误的局部规则，构造反例。能解释错误贪心为何失败，才算理解正确贪心。

\textbf{复杂度变化。} 暴力选择箭位置复杂；贪心排序 O(n log n)，扫描 O(n)，空间取决于排序。

\textbf{迁移追问。} 如果题目改成“它和无重叠区间是同一类最早结束贪心”，先判断原状态是否仍然覆盖新答案集合，再决定是微调边界、改 value 语义，还是换算法。

\noindent\textbf{LeetCode 630 课程表 III。}

\textbf{题目说明。} LeetCode 630 课程表 III 的题面入口要先还原成具体任务：每门课有持续时间和截止日，已选课程总时长不能超过当前截止日。

\textbf{样例入口。} 拿课程 [100, 200]、[200, 1300]、[1000, 1250] 手算，按截止日排序后若总时长超出当前截止日，就丢掉已选课程里耗时最长的。

\textbf{暴力基线。} 暴力枚举课程子集和排列，检查是否能在各自截止日前完成，取最多课程数。

\textbf{暴力过程。} 课程 [100, 200]、[200, 1300]、[1000, 1250] 中，按截止日处理时，若加入一门课后总时长超出当前截止日，就要考虑删掉已选中耗时最长的课。

\textbf{重复工作。} 题面模型：每门课有持续时间和截止日，已选课程总时长不能超过当前截止日。暴力版的慢点要落到本题：浪费在枚举所有选择。按截止日排序后，当前前缀课程只需要让总耗时尽量小；超时删除最长课程给未来留下最多时间。后面的优化只消掉这类重复工作，不能改变答案集合。

\textbf{优化条件。} 本题的关键观察是：按截止日排序，先选当前课；若超时，就用大顶堆丢掉已选中耗时最长的课。这句话必须能翻译成可维护状态；如果题目条件改变后观察不再成立，就不能继续套原来的写法。

\textbf{相邻状态。} 规律是大顶堆保存已选课程时长，替换最长课能给未来留下最多时间。把这个特例继续拆开看：先构造一个非贪心选择，再比较两者留下的资源、右边界或剩余时间。只有能把非贪心第一步交换成贪心第一步且不变差，局部选择才是规律。

\textbf{状态复用。} 按 deadline 升序排序。扫描课程时先加入 duration 到 maxHeap 并累加 total\_time；若 total\_time>deadline，就弹出堆中最长 duration。手算时只改被新输入影响的那一小部分，而不是回头重算完整候选。

\textbf{指针和字段。} 状态字段要能说清：total\_time 是已选课程总耗时，maxHeap 保存已选课程时长，deadline 是当前必须满足的最紧约束。手算时按这个协议跟踪：加入 1000 后若超过截止日，从已选课程中丢掉耗时最长的一门，课程数量尽量保留，且总耗时降得最多。

\textbf{代码落点。} 关键是大顶堆保存 duration，不保存收益。每次超时只弹一次最长课，因为刚加入前状态已经满足之前截止日。关键代码行必须能逐句翻译成“旧状态怎样变成新状态”，否则就只是模板。

\textbf{正确性。} 在相同课程数量下，总耗时越短越不妨碍后续课程；超时后删除最长课程能得到同数量减一情况下最短总耗时。

\textbf{边界实验。} 先用课程无法单独完成、截止日相同、替换已选课程、空列表做反例检查。实验时除了课程无法单独完成、截止日相同、替换已选课程、空列表，还要故意替换成一个看似合理但错误的局部规则，构造反例。能解释错误贪心为何失败，才算理解正确贪心。

\textbf{复杂度变化。} 暴力指数级；排序 O(n log n)，每门课入堆一次、可能出堆一次，O(n log n) 时间，O(n) 空间。

\textbf{迁移追问。} 如果题目改成“这是反悔贪心：接受后在约束被破坏时删最差候选”，先判断原状态是否仍然覆盖新答案集合，再决定是微调边界、改 value 语义，还是换算法。

\noindent\textbf{LeetCode 646 最长数对链。}

\textbf{题目说明。} LeetCode 646 最长数对链 的题面入口要先还原成具体任务：每个数对像一个区间，选择链要求前一个右端小于后一个左端。

\textbf{样例入口。} 拿 [[1, 2], [2, 3], [3, 4]] 手算，[1, 2] 不能接 [2, 3]，因为要求前一个右端小于后一个左端；可以接 [3, 4]。

\textbf{暴力基线。} 慢版本枚举选择顺序或用 DP 保留多个可能方案，先确认最优值存在。随后才检查局部选择是否能通过交换或领先性固定下来。放到本题就是：先按“每个数对像一个区间，选择链要求前一个右端小于后一个左端”完整试慢版本；样例里已经能看到“规律是按右端升序贪心选择，右端越小越容易接后续数对”，所以优化前必须先能说清慢版本枚举了哪些候选。

\textbf{暴力过程。} 拿 [[1, 2], [2, 3], [3, 4]] 手算，[1, 2] 不能接 [2, 3]，因为要求前一个右端小于后一个左端；可以接 [3, 4]。

\textbf{重复工作。} 题面模型：每个数对像一个区间，选择链要求前一个右端小于后一个左端。暴力版的慢点要落到本题：慢版本会围绕“每个数对像一个区间，选择链要求前一个右端小于后一个左端”枚举选择顺序。样例规律是“规律是按右端升序贪心选择，右端越小越容易接后续数对”，说明存在一个局部选择或边界状态可以安全保留；不能证明这个选择不变差时，就不能把它叫贪心。后面的优化只消掉这类重复工作，不能改变答案集合。

\textbf{优化条件。} 本题的关键观察是：按右端升序选择，当前右端越小越容易接后续数对。这句话必须能翻译成可维护状态；如果题目条件改变后观察不再成立，就不能继续套原来的写法。

\textbf{相邻状态。} 规律是按右端升序贪心选择，右端越小越容易接后续数对。把这个特例继续拆开看：先构造一个非贪心选择，再比较两者留下的资源、右边界或剩余时间。只有能把非贪心第一步交换成贪心第一步且不变差，局部选择才是规律。

\textbf{状态复用。} 把指数级选择压成一个可证明安全的局部选择；状态通常是当前最远边界、最早结束时间、最小代价或剩余资源。本题落点是：规律是按右端升序贪心选择，右端越小越容易接后续数对。手算时只改被新输入影响的那一小部分，而不是回头重算完整候选。

\textbf{指针和字段。} 状态字段要能说清：处理顺序、当前已选方案、剩余资源或最远边界、答案计数；若使用排序，要说明排序键；每次局部选择后要能说明当前状态不落后。手算时按这个协议跟踪：拿 [[1, 2], [2, 3], [3, 4]] 手算，[1, 2] 不能接 [2, 3]，因为要求前一个右端小于后一个左端；可以接 [3, 4]。然后按协议复查：手算时同时写一个非贪心选择，与贪心选择比较剩余资源。若无法说明贪心后剩余资源不差，就不能直接下结论。

\textbf{代码落点。} 把局部规则写成业务含义；若需要排序，就处理相同关键字；循环中只维护证明需要的状态，不把局部选择和答案更新混在一起。关键代码行必须能逐句翻译成“旧状态怎样变成新状态”，否则就只是模板。

\textbf{正确性。} 证明从交换或领先性开始：任意最优解都能替换成当前贪心选择而不变差，或贪心状态始终不落后。

\textbf{边界实验。} 先用端点相等不允许连接、完全覆盖、负数端点、空数组做反例检查。实验时除了端点相等不允许连接、完全覆盖、负数端点、空数组，还要故意替换成一个看似合理但错误的局部规则，构造反例。能解释错误贪心为何失败，才算理解正确贪心。

\textbf{复杂度变化。} 暴力枚举选择顺序可能指数级；贪心通常先排序，再线性扫描或用堆维护少量候选。

\textbf{迁移追问。} 如果题目改成“也可以 DP，但贪心用最早结束证明更简洁”，先判断原状态是否仍然覆盖新答案集合，再决定是微调边界、改 value 语义，还是换算法。

\noindent\textbf{LeetCode 678 有效的括号字符串。}

\textbf{题目说明。} LeetCode 678 有效的括号字符串 的题面入口要先还原成具体任务：星号可以作为左括号、右括号或空，使未匹配左括号数量变成一个可能区间。

\textbf{样例入口。} 拿 (*) 手算，星号可以当空或左括号；拿 (*))，星号必须当左括号才能匹配后面的右括号。与其枚举星号三种选择，不如维护可能未匹配左括号数的区间 [low, high]。

\textbf{暴力基线。} 慢版本枚举选择顺序或用 DP 保留多个可能方案，先确认最优值存在。随后才检查局部选择是否能通过交换或领先性固定下来。放到本题就是：先按“星号可以作为左括号、右括号或空，使未匹配左括号数量变成一个可能区间”完整试慢版本；样例里已经能看到“规律是左括号让区间加一，右括号让区间减一，星号让下界减一上界加一，且下界不能低于零”，所以优化前必须先能说清慢版本枚举了哪些候选。

\textbf{暴力过程。} 拿 (*) 手算，星号可以当空或左括号；拿 (*))，星号必须当左括号才能匹配后面的右括号。与其枚举星号三种选择，不如维护可能未匹配左括号数的区间 [low, high]。

\textbf{重复工作。} 题面模型：星号可以作为左括号、右括号或空，使未匹配左括号数量变成一个可能区间。暴力版的慢点要落到本题：慢版本会围绕“星号可以作为左括号、右括号或空，使未匹配左括号数量变成一个可能区间”枚举选择顺序。样例规律是“规律是左括号让区间加一，右括号让区间减一，星号让下界减一上界加一，且下界不能低于零”，说明存在一个局部选择或边界状态可以安全保留；不能证明这个选择不变差时，就不能把它叫贪心。后面的优化只消掉这类重复工作，不能改变答案集合。

\textbf{优化条件。} 本题的关键观察是：贪心维护可能未匹配左括号数的下界和上界，遇到星号扩大区间并把下界截到零。这句话必须能翻译成可维护状态；如果题目条件改变后观察不再成立，就不能继续套原来的写法。

\textbf{相邻状态。} 规律是左括号让区间加一，右括号让区间减一，星号让下界减一上界加一，且下界不能低于零。把这个特例继续拆开看：先构造一个非贪心选择，再比较两者留下的资源、右边界或剩余时间。只有能把非贪心第一步交换成贪心第一步且不变差，局部选择才是规律。

\textbf{状态复用。} 把指数级选择压成一个可证明安全的局部选择；状态通常是当前最远边界、最早结束时间、最小代价或剩余资源。本题落点是：规律是左括号让区间加一，右括号让区间减一，星号让下界减一上界加一，且下界不能低于零。手算时只改被新输入影响的那一小部分，而不是回头重算完整候选。

\textbf{指针和字段。} 状态字段要能说清：处理顺序、当前已选方案、剩余资源或最远边界、答案计数；若使用排序，要说明排序键；每次局部选择后要能说明当前状态不落后。手算时按这个协议跟踪：拿 (*) 手算，星号可以当空或左括号；拿 (*))，星号必须当左括号才能匹配后面的右括号。与其枚举星号三种选择，不如维护可能未匹配左括号数的区间 [low, high]。然后按协议复查：手算时同时写一个非贪心选择，与贪心选择比较剩余资源。若无法说明贪心后剩余资源不差，就不能直接下结论。

\textbf{代码落点。} 把局部规则写成业务含义；若需要排序，就处理相同关键字；循环中只维护证明需要的状态，不把局部选择和答案更新混在一起。关键代码行必须能逐句翻译成“旧状态怎样变成新状态”，否则就只是模板。

\textbf{正确性。} 证明从交换或领先性开始：任意最优解都能替换成当前贪心选择而不变差，或贪心状态始终不落后。

\textbf{边界实验。} 先用上界变负、最终下界不为零、连续星号、前缀右括号过多做反例检查。实验时除了上界变负、最终下界不为零、连续星号、前缀右括号过多，还要故意替换成一个看似合理但错误的局部规则，构造反例。能解释错误贪心为何失败，才算理解正确贪心。

\textbf{复杂度变化。} 暴力枚举选择顺序可能指数级；贪心通常先排序，再线性扫描或用堆维护少量候选。

\textbf{迁移追问。} 如果题目改成“也可用双栈保存左括号和星号位置，但区间贪心更能说明状态压缩”，先判断原状态是否仍然覆盖新答案集合，再决定是微调边界、改 value 语义，还是换算法。

\noindent\textbf{LeetCode 763 划分字母区间。}

\textbf{题目说明。} LeetCode 763 划分字母区间 的题面入口要先还原成具体任务：每个片段必须覆盖其中所有字符的最后出现位置。

\textbf{样例入口。} 拿 ababcbacadefegdehijhklij 手算，字符 a 的最后位置在 8，所以第一个片段至少到 8；扫描中 b、c 的最后位置也不超过 8，到 8 才能切。

\textbf{暴力基线。} 暴力尝试每个切分位置，并检查切出来的片段中是否有字符还会在后面再次出现。

\textbf{暴力过程。} 字符串 ababcbacadefegdehijhklij 中，a 最后出现在 8，所以第一个片段不可能早于 8 结束；扫描到 b、c 时也要把片段右端扩到它们的最后位置。

\textbf{重复工作。} 题面模型：每个片段必须覆盖其中所有字符的最后出现位置。暴力版的慢点要落到本题：浪费在每次试切后重新检查片段字符。每个字符的最后出现位置可以预处理，扫描时维护当前片段必须覆盖的最远右端。后面的优化只消掉这类重复工作，不能改变答案集合。

\textbf{优化条件。} 本题的关键观察是：扫描维护当前片段最远右端，到达右端即可切分。这句话必须能翻译成可维护状态；如果题目条件改变后观察不再成立，就不能继续套原来的写法。

\textbf{相邻状态。} 规律是当前片段右端是片段内所有字符最后出现位置的最大值。把这个特例继续拆开看：先构造一个非贪心选择，再比较两者留下的资源、右边界或剩余时间。只有能把非贪心第一步交换成贪心第一步且不变差，局部选择才是规律。

\textbf{状态复用。} last[ch] 保存字符最后出现下标。扫描字符串时，right=max(right, last[s[i]])；当 i==right，当前片段可以切开。手算时只改被新输入影响的那一小部分，而不是回头重算完整候选。

\textbf{指针和字段。} 状态字段要能说清：start 是当前片段起点，i 是扫描位置，right 是当前片段内所有字符最后位置的最大值。手算时按这个协议跟踪：第一个片段从 0 开始，right 先因 a 变成 8；扫描到 i=8 时，片段内所有字符都不会再出现，可以切出长度 9。

\textbf{代码落点。} 关键是先预处理 last，再扫描更新 right。切分后记录 right-start+1，并令 start=i+1。关键代码行必须能逐句翻译成“旧状态怎样变成新状态”，否则就只是模板。

\textbf{正确性。} 片段必须覆盖内部每个字符的最后出现位置，所以 right 之前不能切；当 i==right 时，片段内字符都不会越界，立即切能最大化片段数量。

\textbf{边界实验。} 先用单字符、所有字符相同、字符多次跨段做反例检查。实验时除了单字符、所有字符相同、字符多次跨段，还要故意替换成一个看似合理但错误的局部规则，构造反例。能解释错误贪心为何失败，才算理解正确贪心。

\textbf{复杂度变化。} 暴力试切最坏 O(\ensuremath{n^{2}})，预处理加扫描 O(n) 时间，字符集空间 O(1) 或 O(字符集)。

\textbf{迁移追问。} 如果题目改成“这是最早合法切分贪心，能最大化片段数量”，先判断原状态是否仍然覆盖新答案集合，再决定是微调边界、改 value 语义，还是换算法。

\noindent\textbf{LeetCode 860 柠檬水找零。}

\textbf{题目说明。} LeetCode 860 柠檬水找零 的题面入口要先还原成具体任务：找零资源只有 5 元和 10 元，5 元比 10 元更灵活。

\textbf{样例入口。} 拿 bills=[5, 5, 5, 10, 20] 手算，前三个人给 5 后手里有三个 5；收到 10 找一个 5；收到 20 时优先找 10+5，保留更多 5 给后续 10 找零。

\textbf{暴力基线。} 慢版本枚举选择顺序或用 DP 保留多个可能方案，先确认最优值存在。随后才检查局部选择是否能通过交换或领先性固定下来。放到本题就是：先按“找零资源只有 5 元和 10 元，5 元比 10 元更灵活”完整试慢版本；样例里已经能看到“规律是维护 5 元和 10 元数量，20 找零时优先使用较不灵活的 10 元”，所以优化前必须先能说清慢版本枚举了哪些候选。

\textbf{暴力过程。} 拿 bills=[5, 5, 5, 10, 20] 手算，前三个人给 5 后手里有三个 5；收到 10 找一个 5；收到 20 时优先找 10+5，保留更多 5 给后续 10 找零。

\textbf{重复工作。} 题面模型：找零资源只有 5 元和 10 元，5 元比 10 元更灵活。暴力版的慢点要落到本题：慢版本会围绕“找零资源只有 5 元和 10 元，5 元比 10 元更灵活”枚举选择顺序。样例规律是“规律是维护 5 元和 10 元数量，20 找零时优先使用较不灵活的 10 元”，说明存在一个局部选择或边界状态可以安全保留；不能证明这个选择不变差时，就不能把它叫贪心。后面的优化只消掉这类重复工作，不能改变答案集合。

\textbf{优化条件。} 本题的关键观察是：收到 10 用一个 5 找零；收到 20 优先用 10+5，不能时再用三个 5。这句话必须能翻译成可维护状态；如果题目条件改变后观察不再成立，就不能继续套原来的写法。

\textbf{相邻状态。} 规律是维护 5 元和 10 元数量，20 找零时优先使用较不灵活的 10 元。把这个特例继续拆开看：先构造一个非贪心选择，再比较两者留下的资源、右边界或剩余时间。只有能把非贪心第一步交换成贪心第一步且不变差，局部选择才是规律。

\textbf{状态复用。} 把指数级选择压成一个可证明安全的局部选择；状态通常是当前最远边界、最早结束时间、最小代价或剩余资源。本题落点是：规律是维护 5 元和 10 元数量，20 找零时优先使用较不灵活的 10 元。手算时只改被新输入影响的那一小部分，而不是回头重算完整候选。

\textbf{指针和字段。} 状态字段要能说清：处理顺序、当前已选方案、剩余资源或最远边界、答案计数；若使用排序，要说明排序键；每次局部选择后要能说明当前状态不落后。手算时按这个协议跟踪：拿 bills=[5, 5, 5, 10, 20] 手算，前三个人给 5 后手里有三个 5；收到 10 找一个 5；收到 20 时优先找 10+5，保留更多 5 给后续 10 找零。然后按协议复查：手算时同时写一个非贪心选择，与贪心选择比较剩余资源。若无法说明贪心后剩余资源不差，就不能直接下结论。

\textbf{代码落点。} 把局部规则写成业务含义；若需要排序，就处理相同关键字；循环中只维护证明需要的状态，不把局部选择和答案更新混在一起。关键代码行必须能逐句翻译成“旧状态怎样变成新状态”，否则就只是模板。

\textbf{正确性。} 证明从交换或领先性开始：任意最优解都能替换成当前贪心选择而不变差，或贪心状态始终不落后。

\textbf{边界实验。} 先用第一张不是 5、连续 10、连续 20、只有 5、5 元数量不足做反例检查。实验时除了第一张不是 5、连续 10、连续 20、只有 5、5 元数量不足，还要故意替换成一个看似合理但错误的局部规则，构造反例。能解释错误贪心为何失败，才算理解正确贪心。

\textbf{复杂度变化。} 暴力枚举选择顺序可能指数级；贪心通常先排序，再线性扫描或用堆维护少量候选。

\textbf{迁移追问。} 如果题目改成“若面额集合变化，需要重新证明哪种资源更灵活”，先判断原状态是否仍然覆盖新答案集合，再决定是微调边界、改 value 语义，还是换算法。

\subsection{实验复盘}

追问通常会要求证明。请准备交换论证或领先性论证，并给出一个看似合理但错误的贪心规则反例。

复盘本节时先从 LeetCode 45 跳跃游戏 II、LeetCode 55 跳跃游戏、LeetCode 122 买卖股票 II 开始：每道题都先手写一个能覆盖全部候选的暴力版，再指出优化结构具体保存了哪一段历史信息，最后用相邻案例比较为什么不能互换解法。

